\chapter{Ground states of gapped chains in polynomial time}
\label{ch:gapped_chain_algorithms}

Landau, Vazirani and Vidick~\cite{Landau2015Polynomial} showed that a
matrix product state approximating the ground state of a gapped chain can
be computed in polynomial time, and Arad, Landau, Vazirani and
Vidick~\cite{Arad2016RigorousRG} gave a more efficient renormalization
algorithm. Both build the approximation one cut at a time. At each cut they
keep a small set of vectors on the sites to the left of the cut, a viable
set, whose span still contains the left half of a good approximation. This
chapter formalizes viable sets, trimming them to polynomial size, gluing a
left state at a cut, and the polynomial filters with their families of
coefficients on two windows.

\section{Viable sets for the ground state of a gapped chain}
\label{sec:lvv_viable_sets}

Landau, Vazirani and Vidick build an approximation to the ground state
$\Gamma$ of a gapped chain of $L$ sites cut by cut, keeping at the cut after
$i$ sites a small set $S$ of vectors on the first $i$ sites that supports the
left half of a good approximation to
$\Gamma$~\cite[lines 240--244]{Landau2015Polynomial}. For a vector $\psi$ on
the chain, write $M_i(\psi)$ for its coefficient matrix across the cut after
$i$ sites, with rows indexed by the configurations of the first $i$ sites.
The reduced state of $\psi$ on the first $i$ sites is $M_i(\psi)M_i(\psi)^*$,
whose range is the span of the columns of $M_i(\psi)$.

\begin{definition}[Viable sets]
    \label{def:lvv_viable_set}
    \lean{TNLean.MPS.AreaLaw.cutColumn,
        TNLean.MPS.AreaLaw.leftSupportSubmodule,
        TNLean.MPS.AreaLaw.mem_leftSupportSubmodule,
        TNLean.MPS.AreaLaw.IsViable}
    \leanok
    \uses{def:chain_model}
    Let $\Gamma$ be a vector on the chain, $\delta\in\mathbb R$ and $S$ a set of
    vectors on the first $i$ sites. Then $S$ is $(i,\delta)$-viable for $\Gamma$
    if some unit vector $\phi$ with $|\langle\phi,\Gamma\rangle|\ge1-\delta$ has
    every column of $M_i(\phi)$ in $\operatorname{Span}(S)$, that is, $\phi$ lies
    in $\operatorname{Span}(S)\otimes\mathcal H_R$ and its reduced state on the
    first $i$ sites is supported on $\operatorname{Span}(S)$. This is
    Definition~def:viable of~\cite[lines 240--242]{Landau2015Polynomial}.
\end{definition}

\begin{theorem}[Support through the projection onto the span]
    \label{thm:lvv_viable_projection}
    \lean{TNLean.MPS.AreaLaw.spanProjection,
        TNLean.MPS.AreaLaw.isStarProjection_spanProjection,
        TNLean.MPS.AreaLaw.isHermitian_spanProjection,
        TNLean.MPS.AreaLaw.spanProjection_mul_self,
        TNLean.MPS.AreaLaw.spanProjection_mul_cutMatrix_eq_iff,
        TNLean.MPS.AreaLaw.regionState_leftBlock_eq,
        TNLean.MPS.AreaLaw.spanProjection_mul_regionState_eq_iff,
        TNLean.MPS.AreaLaw.leftApply_cutVector_eq_iff,
        TNLean.MPS.AreaLaw.leftApply_spanProjection_cutVector_eq_iff}
    \leanok
    \uses{def:lvv_viable_set}
    Let $P_S$ be the orthogonal projection onto $\operatorname{Span}(S)$, a
    Hermitian idempotent matrix on the configurations of the first $i$ sites.
    For every vector $\psi$ the following are equivalent: the columns of
    $M_i(\psi)$ lie in $\operatorname{Span}(S)$; $P_SM_i(\psi)=M_i(\psi)$;
    $P_S\rho=\rho$ for the reduced state $\rho=M_i(\psi)M_i(\psi)^*$ of $\psi$ on
    the first $i$ sites.
\end{theorem}

\begin{proof}
    \leanok
    A column $c$ of $M_i(\psi)$ lies in $\operatorname{Span}(S)$ exactly when
    $P_Sc=c$. The range of $M_i(\psi)M_i(\psi)^*$ is the column space of
    $M_i(\psi)$.
\end{proof}

\begin{theorem}[Viable sets and viable subspaces]
    \label{thm:lvv_viable_closeness}
    \lean{TNLean.MPS.AreaLaw.IsViable.le_norm_starProjection,
        TNLean.MPS.AreaLaw.isViable_iff_le_norm_starProjection,
        TNLean.MPS.AreaLaw.IsViable.forall_mem_span_singleton,
        TNLean.MPS.AreaLaw.isViable_iff_forall_mem_span_singleton}
    \leanok
    \uses{def:lvv_viable_set}
    Let $U=\operatorname{Span}(S)\otimes\mathcal H_R$ and let $P_U$ be the
    orthogonal projection onto $U$. If $S$ is $(i,\delta)$-viable for $\Gamma$,
    then $\|P_U\Gamma\|\ge1-\delta$, and for $\delta<1$ the converse holds. For a
    unit vector $\Gamma$ and $\delta\le1$, viability gives
    $(1-\delta')\|t\|^2\le\|P_Ut\|^2$ for every $t$ in the span of $\Gamma$, with
    $\delta'=1-(1-\delta)^2$, and for $\delta<1$ the converse holds. This is the
    closeness of $U$ to the span of $\Gamma$ in the sense of the viable subspaces
    of Arad, Landau, Vazirani and
    Vidick~\cite[lines 433--440]{Arad2016RigorousRG}, with $\delta'$ in place of
    $\delta$, since $\delta$ bounds an overlap and $\delta'$ a squared overlap.
\end{theorem}

\begin{proof}
    \leanok
    For a unit $\phi\in U$, $|\langle\phi,\Gamma\rangle|
        =|\langle\phi,P_U\Gamma\rangle|\le\|P_U\Gamma\|$, with equality for
    $\phi=P_U\Gamma/\|P_U\Gamma\|$ when $P_U\Gamma\ne0$. Squaring gives the second
    form.
\end{proof}

\begin{definition}[Viable sets of matrix product states]
    \label{def:lvv_viable_mps}
    \lean{TNLean.MPS.AreaLaw.leftConfigOfFin,
        TNLean.MPS.AreaLaw.IsLeftOpenMPS,
        TNLean.MPS.AreaLaw.IsViableMPS}
    \leanok
    \uses{def:lvv_viable_set, def:chain_open_mps}
    The set $S$ is $(i,s,b,\delta)$-viable for $\Gamma$ if it is
    $(i,\delta)$-viable, has at most $s$ members, and each member is an open
    matrix product state on the first $i$ sites of bond dimension at most
    $\max(1,b)$~\cite[line 243]{Landau2015Polynomial}. Bond dimension $0$ is read
    as $1$; the source's bond dimensions are
    positive~\cite{gap:lvv13_bond_dimension_zero}.
\end{definition}

\begin{theorem}[Extending a viable set by one site]
    \label{thm:lvv_viable_extension}
    \lean{TNLean.MPS.AreaLaw.IsOpenMPS.snoc,
        TNLean.MPS.AreaLaw.restrictLeft,
        TNLean.MPS.AreaLaw.extendSite,
        TNLean.MPS.AreaLaw.lastLeftSite,
        TNLean.MPS.AreaLaw.extendSite_apply,
        TNLean.MPS.AreaLaw.extend,
        TNLean.MPS.AreaLaw.extendSite_mem_extend,
        TNLean.MPS.AreaLaw.card_extend_le,
        TNLean.MPS.AreaLaw.consRightConfig,
        TNLean.MPS.AreaLaw.cutColumn_eq_sum_extendSite,
        TNLean.MPS.AreaLaw.leftSupportSubmodule_le_extend,
        TNLean.MPS.AreaLaw.IsViable.extend,
        TNLean.MPS.AreaLaw.IsLeftOpenMPS.extendSite,
        TNLean.MPS.AreaLaw.IsViableMPS.extend}
    \leanok
    \uses{def:lvv_viable_mps}
    Let $1\le i\le L$, let $d$ be the local dimension, and let $S'$ be the set
    of vectors $v\otimes|k\rangle$ on the first $i$ sites, for $v\in S$ on the
    first $i-1$ sites and $k$ a basis state. If $S$ is
    $(i-1,s,b,\delta)$-viable for $\Gamma$, then $S'$ is
    $(i,ds,b,\delta)$-viable for $\Gamma$, with the same witness. This is
    Claim~c:ext of~\cite[lines 379--400]{Landau2015Polynomial}, stated there
    for $\delta=c_\varepsilon/n$.
\end{theorem}

\begin{proof}
    \leanok
    \uses{def:lvv_viable_mps}
    $S'$ has at most $d|S|$ members. Appending a basis state to an open matrix
    product state adds a bond of rank one, so the bond dimension stays at most
    $\max(1,b)$. A column of $M_i(\phi)$ is a sum over the state of site $i$ of
    vectors $c\otimes|k\rangle$ with $c$ a column of $M_{i-1}(\phi)$, and these
    lie in $\operatorname{Span}(S')$ when the columns of $M_{i-1}(\phi)$ lie in
    $\operatorname{Span}(S)$; the witness $\phi$ is unchanged.
\end{proof}

\section{Trimming across one cut}

Landau, Vazirani and Vidick~\cite{Landau2015Polynomial} (arXiv:1307.5143v1,
lines 330--333) trim a chain vector across a cut by keeping the Schmidt terms
with the largest coefficients. Trimming across one cut does not increase the
Schmidt rank across any other cut (lines 544--553).

\begin{definition}[Trimming a chain vector]
    \label{def:chain_trim}
    \lean{TNLean.MPS.AreaLaw.chainTrim}
    \leanok
    \uses{def:chain_open_mps}
    Let $\psi$ be a vector on $L$ sites of local dimension $q$, let $i$ and
    $m$ be natural numbers, and write $C=M_i(\psi)$ for the coefficient matrix
    at the cut $i$. A \emph{trimming projection} of $\psi$ at $m$ across the
    cut $i$ is an orthogonal projection $P$ on $\mathcal H_{[0,i)}$ of rank at
    most $m$ with $\operatorname{tr}(PCC^{*})$ equal to the sum of the $m$
    largest eigenvalues of $CC^{*}$. The \emph{trimmed vector} is
    $\operatorname{trim}^i_P\psi=(P\otimes1)\psi$, whose coefficient matrix
    at the cut $i$ is $PC$. When a Schmidt coefficient is repeated, every
    choice of kept eigenvectors gives a trimming projection.
\end{definition}

\begin{theorem}[Trimming does not increase cut ranks]
    \label{thm:chain_trim_cut_ranks}
    \lean{
        TNLean.MPS.AreaLaw.cutMatrix_chainTrim,
        TNLean.MPS.AreaLaw.cutVector_chainTrim,
        TNLean.MPS.AreaLaw.rank_cutMatrix_chainTrim_self_le,
        TNLean.MPS.AreaLaw.rank_cutMatrix_chainTrim_le
    }
    \leanok
    \uses{def:chain_trim}
    Let $P$ be a trimming projection of $\psi$ at $m$ across the cut $i$.
    The coefficient matrix of $\operatorname{trim}^i_P\psi$ at the cut $i$ is
    $PM_i(\psi)$, so $\operatorname{rank}M_i(\operatorname{trim}^i_P\psi)\le m$, and
    $\operatorname{rank}M_j(\operatorname{trim}^i_P\psi)\le
        \operatorname{rank}M_j(\psi)$ at every cut $j$. No assumption on $L$
    or $q$ is needed.
\end{theorem}
\begin{proof}\leanok
    \uses{thm:hastings_supported_cut_rank}
    The first bound holds because $PC$ has rank at most that of $P$. For
    $j\ge i$ the operator $P\otimes1$ acts on sites to the left of the cut
    $j$. For $j\le i$, the projection $P$ commutes with $CC^{*}$, since it
    maximizes $\operatorname{tr}(RCC^{*})$ among orthogonal projections $R$
    of the same trace. Hence $PC=CQ^{\mathsf T}$, where $Q^{\mathsf T}$ is the
    support projection of $C^{*}PC$, so $(P\otimes1)\psi=(1\otimes Q)\psi$
    with $Q$ acting on sites to the right of the cut $j$. In both cases the
    operator acts on one side of the cut $j$ and cannot increase its rank.
    If there are no configurations, both coefficient matrices are empty.
\end{proof}

\begin{theorem}[Trimming error on the chain]
    \label{thm:chain_trim_error}
    \lean{
        TNLean.MPS.AreaLaw.norm_chainTrim_le,
        TNLean.MPS.AreaLaw.norm_inner_sub_chainTrim_le
    }
    \leanok
    \uses{def:chain_trim}
    Let $P$ be a trimming projection of $\psi$ at $m$ across the cut $i$,
    and let $w$ be a vector with $\operatorname{rank}M_i(w)\le D$. Then
    $\|\operatorname{trim}^i_P\psi\|\le\|\psi\|$ and
    \begin{align}
        |\langle\psi-\operatorname{trim}^i_P\psi,w\rangle|
        \le\|\psi\|\,\|w\|\sqrt{\frac{D}{m+1}}. \notag
    \end{align}
    The source states error $\epsilon$ for $m=D/\epsilon$, which is false;
    the bound above gives error $\epsilon$ for $m+1\ge D/\epsilon^2$
    \cite{gap:lvv13_trimming_radius}.
\end{theorem}
\begin{proof}\leanok
    The cut coordinates at $i$ preserve norms and inner products, and in
    them the trimmed vector is $(P\otimes1)$ applied to the cut vector of
    $\psi$. The squared norm of $(P\otimes1)\psi$ is the sum of the $m$
    largest eigenvalues of $CC^{*}$, at most their total $\|\psi\|^2$. For the
    second bound, let $T$ be the support projection of the part of $M_i(w)$
    orthogonal to the range of $P$. Then $P+T$ is an orthogonal projection of
    rank at most $m+D$, so by Ky Fan's maximum principle the discarded part
    paired with $w$ captures at most $D$ times the $(m+1)$-st eigenvalue of
    $CC^{*}$, which is at most $\|\psi\|^2/(m+1)$.
\end{proof}

\subsection{Constant rank across one cut}

Landau, Vazirani and Vidick~\cite{Landau2015Polynomial} (lines 342--364)
use that the ground state of a gapped chain can be trimmed across one cut to
a bond dimension independent of the length of the chain, and that trimming a
vector of low energy keeps its energy low. The bond dimension they state,
$\exp(O(\varepsilon^{-1}\log^3d\log(1/\delta)))$, comes from Arad, Kitaev,
Landau and Vazirani~\cite{Arad2013AreaLaw}; the one below comes from the
area law of this chapter and has different constants
\cite{gap:lvv13_constant_bond}. The energy bound after trimming is the
corrected one of \cite{gap:lvv13_trimming_energy}.

\begin{definition}[A constant bond dimension]
    \label{def:chain_constant_bond}
    \lean{TNLean.MPS.AreaLaw.constantCutBond}
    \leanok
    \uses{def:chain_model}
    For a chain of local dimension $q$ with range $R$, term norms at most $J$
    and gap $\Delta>0$, let $C_{\mathrm{int}}$ be the interval entropy bound,
    $s=\sqrt{(1+J/\Delta)B_R}$, and $0<\delta\le1$. Put
    \begin{align}
        B_\delta=\Bigl\lfloor\exp\Bigl(C_{\mathrm{int}}
        +32s\log\frac{2e^{e/2}}{\delta}\Bigr)\Bigr\rfloor. \notag
    \end{align}
    It depends neither on the length of the chain nor on a cut.
\end{definition}

\begin{theorem}[Trimming the ground state to a constant bond dimension]
    \label{thm:chain_constant_bond_overlap}
    \lean{
        TNLean.MPS.AreaLaw.ChainHamiltonian.one_sub_le_re_inner_chainTrim,
        TNLean.MPS.AreaLaw.inner_normalized_chainTrim,
        TNLean.MPS.AreaLaw.cutVector_normalized_chainTrim,
        TNLean.MPS.AreaLaw.regionState_leftBlock_eq_leftMarginal
    }
    \leanok
    \uses{def:chain_trim, def:chain_constant_bond}
    Let $\Omega$ be a gapped ground vector, with gap $\Delta>0$, of a
    finite-range Hamiltonian with term norms at most $J\ge0$ on a chain of
    local dimension $q\ge1$, and let $0<\delta\le1$. For every cut $k$ and
    every trimming projection $P$ of $\Omega$ at $B_\delta$ across the cut $k$,
    the vector $v=\operatorname{trim}^k_P\Omega/
        \|\operatorname{trim}^k_P\Omega\|$ satisfies
    $\operatorname{Re}\langle v,\Omega\rangle\ge1-\delta$, and the overlap is
    real and nonnegative. This is the statement of the source with the bond
    dimension of Definition~\ref{def:chain_constant_bond} in place of the
    stated one \cite{gap:lvv13_constant_bond}.
\end{theorem}
\begin{proof}\leanok
    \uses{thm:chain_prefix_spectral_bounds}
    The reduced state $\rho$ of $\Omega$ on the first $k$ sites is $CC^{*}$
    for the coefficient matrix $C$ at the cut. Take
    $w=32s\log(2e^{e/2}/\delta)$ in
    Theorem~\ref{thm:chain_prefix_spectral_bounds}: the spectral projection
    there has rank at most $B_\delta$ and discards mass at most $\delta$. By
    Ky Fan's maximum principle the sum of the $B_\delta$ largest eigenvalues
    of $\rho$ is at least $1-\delta$. This sum is
    $\|\operatorname{trim}^k_P\Omega\|^2$, and
    $\langle v,\Omega\rangle=\|\operatorname{trim}^k_P\Omega\|\ge
        \sqrt{1-\delta}\ge1-\delta$.
\end{proof}

\begin{theorem}[The energy after trimming]
    \label{thm:chain_trim_energy}
    \lean{
        TNLean.MPS.AreaLaw.re_inner_normalized_chainTrim_le,
        TNLean.MPS.AreaLaw.ChainHamiltonian.re_inner_normalized_chainTrim_le,
        TNLean.MPS.AreaLaw.rank_cutMatrix_normalized_chainTrim_le,
        TNLean.MPS.AreaLaw.norm_normalized_chainTrim_eq_one,
        Matrix.IsGappedGroundState.posSemidef_sub_smul_one
    }
    \leanok
    \uses{def:chain_trim, thm:chain_constant_bond_overlap}
    Let $H=H_L+H_R+K$ on a chain, with $H_L$ supported on the first $k$
    sites, $H_R$ on the remaining sites and $0\le K\le1$, and let $\Omega$
    be a gapped ground vector of $H$ with energy $E_0$ and gap $\Delta>0$.
    Let $B$ be a bond dimension and $P_\Omega$ a trimming projection of
    $\Omega$ at $B$ across the cut $k$ whose normalized trimmed vector $u$
    satisfies $\operatorname{Re}\langle u,\Omega\rangle\ge1-\delta_B$. Let
    $w$ be a unit vector with $\operatorname{Re}\langle w,Hw\rangle\le
        E_0+\delta$, and $P$ a trimming projection of $w$ at $B$ across the cut
    $k$. If $\eta\ge2(\delta/\Delta+\delta_B)$ and $2\eta<1$, the normalized
    trimmed vector $v$ of $w$ satisfies
    \begin{align}
        \operatorname{Re}\langle v,Hv\rangle
        \le E_0+\frac{\delta+2\sqrt{2\eta}}{1-2\eta}. \notag
    \end{align}
    For a finite-range chain Hamiltonian that splits in this way at the cut
    $k$, this holds at $B=B_{\delta_B}$ for every $0<\delta_B\le1$. In the
    source $H$ is a sum of nearest-neighbour terms $0\le H_i\le1$ and $K$ is
    the term across the cut; its printed bound $E_0+6\sqrt\delta$ does not
    follow from its argument \cite{gap:lvv13_trimming_energy}.
\end{theorem}
\begin{proof}\leanok
    The gap gives $H\ge E_0$, so $\delta\ge0$, and
    $|\langle\Omega,w\rangle|\ge1-\delta/\Delta$. Since $u$, $w$ and $\Omega$
    are unit vectors, the overlap estimate for three unit vectors gives
    $|\langle u,w\rangle|\ge1-2(\delta_B+\delta/\Delta)\ge1-\eta$. In the
    coordinates of the cut $k$, which preserve norms and inner products,
    $H_L$ becomes $A\otimes1$, $H_R$ becomes $1\otimes B'$, and $u$ has
    Schmidt rank at most $B$. The bipartite energy estimate for trimmed
    vectors then gives the bound: the trimming projection commutes with the
    left marginal of $w$, so $A\otimes1+1\otimes B'$ has no cross terms
    between the trimmed part and the discarded part, the cross term of $K$ is
    at most twice the norm of the discarded part, and by the Eckart--Young
    theorem the squared norm of the discarded part is at most $2\eta$.
\end{proof}

\begin{theorem}[A state of constant rank across one cut]
    \label{thm:chain_constant_rank_state}
    \lean{TNLean.MPS.AreaLaw.ChainHamiltonian.exists_rank_cutMatrix_le_constantCutBond}
    \leanok
    \uses{def:chain_constant_bond}
    Let $\Omega$ be a gapped ground vector, with energy $E_0$ and gap
    $\Delta>0$, of a finite-range chain Hamiltonian with term norms at most
    $J\ge0$ and local dimension $q\ge1$, which splits at the cut $k$ as in
    Theorem~\ref{thm:chain_trim_energy}, and let $0<\delta<1/4$. There is a
    unit vector $v$ with $\operatorname{rank}M_k(v)\le B_\delta$,
    $\operatorname{Re}\langle v,\Omega\rangle\ge1-\delta$, and
    \begin{align}
        \operatorname{Re}\langle v,Hv\rangle\le
        E_0+\frac{4\sqrt\delta}{1-4\delta}. \notag
    \end{align}
    The source states the energy bound $E_0+6\sqrt\delta$
    \cite{gap:lvv13_trimming_energy}.
\end{theorem}
\begin{proof}\leanok
    \uses{thm:chain_constant_bond_overlap, thm:chain_trim_energy}
    Take $v$ the normalized trimmed vector of $\Omega$ at $B_\delta$. Its
    overlap is Theorem~\ref{thm:chain_constant_bond_overlap}, which also
    shows that the trimmed vector is nonzero, and its rank across the cut is
    at most $B_\delta$. Theorem~\ref{thm:chain_trim_energy} with $w=\Omega$,
    energy excess $0$, $\delta_B=\delta$ and $\eta=2\delta$ gives
    $E_0+2\sqrt{4\delta}/(1-4\delta)$.
\end{proof}

\section{Collective trimming of a viable subspace}
\label{sec:alvv_collective_trimming}

Arad, Landau, Vazirani and Vidick reduce the Schmidt ranks of a viable
subspace by trimming a basis of it collectively, cut by
cut~\cite[lines 642--697]{Arad2016RigorousRG}. Let $S$ be a subspace of
vectors on the first $i$ sites of a chain of $L$ sites and $T$ a subspace of
vectors on the chain. Write $\mathcal H_A$ for the first $i$ sites and
$\mathcal H_R$ for the remaining $L-i$, so that the cuts $(j:j+1)$,
$1\le j\le i-1$, of the source are the cuts of the chain after $j$ sites.

\begin{definition}[Trimmed subspace]
    \label{def:alvv_collective_trim}
    \lean{TNLean.MPS.AreaLaw.blockCut,
        TNLean.MPS.AreaLaw.chainTrimVector,
        TNLean.MPS.AreaLaw.chainTrimSubspace,
        TNLean.MPS.AreaLaw.chainTrimSubspace_eq_span}
    \leanok
    \uses{def:lvv_viable_set}
    Let $u_1,\dots,u_s$ be an orthonormal basis of $S$ and $\xi>0$, and set
    $\psi=\sum_k u_k\otimes|k\rangle\in\mathcal H_A\otimes\mathbb C^s$. For
    $j=1,\dots,i-1$ let $P^j$ be the orthogonal projection onto the left
    Schmidt vectors, across the cut after $j$ sites, of
    $(P^1\otimes1)\cdots(P^{j-1}\otimes1)\psi$ with Schmidt coefficient at
    least $\xi$. The trimmed vectors are
    $(P^1\otimes1)\cdots(P^{i-1}\otimes1)u_k$, and the trimmed subspace
    $\operatorname{Trim}_\xi(S)$ is their span. This is
    Definition~def:trimming of~\cite[lines 642--656]{Arad2016RigorousRG}.
\end{definition}

\begin{theorem}[Viability in block coordinates]
    \label{thm:alvv_collective_trim_coordinates}
    \lean{TNLean.MPS.AreaLaw.blockSplit,
        TNLean.MPS.AreaLaw.blockSplit_apply,
        TNLean.MPS.AreaLaw.leftSupportSubmodule_eq_comap,
        TNLean.MPS.AreaLaw.isCloseTo_leftSupportSubmodule_iff,
        TNLean.MPS.AreaLaw.blockCut_mono,
        TNLean.MPS.AreaLaw.mem_blockCut,
        TNLean.MPS.AreaLaw.blockCutSubtypeEquiv,
        TNLean.MPS.AreaLaw.blockRestSubtypeEquiv,
        TNLean.MPS.AreaLaw.rank_cutEnvVector_blockSplit}
    \leanok
    \uses{def:alvv_collective_trim}
    Write a vector $\psi$ on the chain in the coordinates
    $\mathcal H_A\otimes\mathcal H_R$. Then $\operatorname{Span}(S)\otimes
        \mathcal H_R$ is $\delta$-close to $T$ exactly when $S$ is $\delta$-viable
    for $T$ in these coordinates. For $1\le j\le i$, the Schmidt rank of $\psi$
    across the first $j$ sites of $\mathcal H_A$, with $\mathcal H_R$ on the
    right, equals the rank of the coefficient matrix $M_j(\psi)$ of the chain
    cut after $j$ sites.
\end{theorem}

\begin{proof}
    \leanok
    The coordinates are a relabelling of the configurations, so they preserve
    inner products, and $\psi$ lies in $\operatorname{Span}(S)\otimes
        \mathcal H_R$ exactly when every column of $M_i(\psi)$ lies in $S$. Both
    bipartitions in the second claim put the first $j$ sites on the left and
    the rest on the right, so the two coefficient matrices differ only by a
    relabelling of rows and columns.
\end{proof}

\begin{theorem}[Trimming a viable subspace on the chain]
    \label{thm:alvv_collective_trim}
    \lean{TNLean.MPS.AreaLaw.isCloseTo_leftSupportSubmodule_chainTrimSubspace,
        TNLean.MPS.AreaLaw.schmidtRank_cutVector_chainTrimVector_le}
    \leanok
    \uses{def:alvv_collective_trim, thm:alvv_collective_trim_coordinates}
    Let $S$ have dimension $s$ and let $\operatorname{Span}(S)\otimes
        \mathcal H_R$ be $\delta$-close to $T$. Suppose every vector of $T$ has
    $\operatorname{rank}M_j\le b$ at every cut $j=1,\dots,i-1$. Then for
    $\xi>0$ the subspace $\operatorname{Trim}_\xi(S)\otimes\mathcal H_R$ is
    $\big(\delta+2\xi\sqrt{(i-1)b}\big)$-close to $T$. Each trimmed vector
    has Schmidt rank at most $s/\xi^2$ across every cut $(j:j+1)$ of
    $\mathcal H_A$, $1\le j\le i-1$.
    This is Lemma~lem:trimming of~\cite[lines 660--697]{Arad2016RigorousRG},
    without the running time. The source states $\delta+\sqrt{ibs}\,\xi$, which
    its argument does not give; the bound above is the one it
    gives~\cite{gap:alvv16_trimming_constant}.
\end{theorem}

\begin{proof}
    \leanok
    In the coordinates of the previous theorem this is the trimming lemma for
    a subspace of a block of sites with the remaining sites as environment.
    The part of a vector of $\operatorname{Span}(S)\otimes\mathcal H_R$
    removed at the cut after $j$ sites pairs with a vector $t$ of $T$ to at
    most $\xi\sqrt b$ times the norm of the piece of $t$ removed there. The
    pieces removed at different cuts are orthogonal, so their norms sum to at
    most $\sqrt{i-1}\,\|t\|$. The rank bound holds because $\psi$ has at most
    $s/\xi^2$ Schmidt coefficients at least $\xi$ across each cut.
\end{proof}

\section{Gluing a left state at a cut of the chain}
\label{sec:chain_gluing}

Fix the cut after $k$ sites of a chain. Splitting off the last site left of
the cut, a configuration has coordinates $((x,y),z)$, with $x$ on the first
$k-1$ sites, $y$ on site $k-1$ and $z$ on the sites right of the cut. This is
the setting of the boundary contractions and the gluing lemma of Landau,
Vazirani and Vidick~\cite[lines 409--447]{Landau2015Polynomial}, proved in the
companion quantum-channel
volume~\cite[``Gluing'', ``The leading eigenvector of a near-optimal left state'',
    ``Three-way cut coordinates'']{QICLean2026Blueprint}.

\begin{definition}[Cut coordinates at a cut of the chain]
    \label{def:chain_gluing_cut}
    \lean{TNLean.MPS.AreaLaw.leftBlock_sub_one_subset, TNLean.MPS.AreaLaw.CutLeft,
        TNLean.MPS.AreaLaw.CutSite, TNLean.MPS.AreaLaw.CutRight,
        TNLean.MPS.AreaLaw.gluingCut, TNLean.MPS.AreaLaw.leftCutOperator,
        TNLean.MPS.AreaLaw.rightCutOperator}
    \leanok
    The coordinates $((x,y),z)$ of a configuration of the chain at the cut
    after $k$ sites. For an operator $H_L$ supported on the first $k$ sites,
    $H_L'$ is its restriction to them; for an operator $K$ supported off the
    first $k-1$ sites, $K'$ is its restriction to site $k-1$ and the sites right
    of the cut.
\end{definition}

\begin{theorem}[The Hamiltonian at a cut]
    \label{thm:chain_gluing_split}
    \lean{TNLean.MPS.AreaLaw.reindex_gluingCut_eq}
    \leanok
    \uses{def:chain_gluing_cut}
    If $H_L$ is supported on the first $k$ sites and $K$ off the first $k-1$
    sites, then in the cut coordinates $H_L+K=H_L'\otimes1+1\otimes K'$.
\end{theorem}

\begin{proof}
    \leanok
    An operator supported on a set of sites acts as the identity on the
    configuration of the other sites. Thus, in the cut coordinates, with
    the reassociation $((x,y),z)\leftrightarrow(x,(y,z))$ for $K$,
    \begin{align}
        H_L & =H_L'\otimes1, \\
        K   & =1\otimes K'.
    \end{align}
    Adding gives the asserted decomposition of $H_L+K$.
\end{proof}

\begin{theorem}[The leading eigenvector at a cut of the chain]
    \label{thm:chain_leading_eigenvector}
    \lean{TNLean.MPS.AreaLaw.exists_leading_eigenvector_chain}
    \leanok
    \uses{def:chain_gluing_cut}
    Let $H=H_L+K$ with $H_L$ supported on the first $k$ sites and $K$ off the
    first $k-1$ sites have a gapped ground state $\Gamma$ with energy $E_0$ and
    gap $\Delta$. Let $W$ be an isometry from a bond $\mathbb C^m$ into the
    sites right of the cut, and let $w=(1\otimes W)x$ in the cut coordinates
    have energy at most $E_0+e$. Let $\sigma$ be a density matrix on the first
    $k$ sites and the bond with
    $\operatorname{Re}\operatorname{tr}(\sigma(H_L'\otimes1))\le
        \operatorname{Re}\operatorname{tr}(\lvert x\rangle\langle x\rvert(H_L'\otimes1))$
    and
    $\lVert\operatorname{tr}_{[1,k-1]}\sigma-\operatorname{tr}_{[1,k-1]}
        \lvert x\rangle\langle x\rvert\rVert_1\le\tau$. If
    $4(e+\lVert K\rVert\tau)<\Delta$, then $\sigma$ has a unit eigenvector $u$
    with eigenvalue larger than $\tfrac12$ such that $(1\otimes W)u$, read on
    the chain, has energy at most $E_0+2(e+\lVert K\rVert\tau)$ and overlap at
    least $1-2(e+\lVert K\rVert\tau)/\Delta$ with $\Gamma$.

    This is Lemma~l:leading of~\cite[lines 493--505]{Landau2015Polynomial} on
    the chain, with $H_L=H_1+\cdots+H_{i-1}$ and $K=H_i+H_R$; the energy excess
    $e$ of the trimmed state of Lemma~l:within is corrected
    in~\cite{gap:lvv13_trimming_energy}.
\end{theorem}

\begin{proof}
    \leanok
    \uses{thm:chain_gluing_split}
    In the cut coordinates the Hamiltonian, the ground state and $w$ satisfy
    the hypotheses of the leading-eigenvector theorem of the companion volume,
    with $\lVert K'\rVert\le\lVert K\rVert$ since $K'$ is a corner of
    $1\otimes K'$. Let $E$ be the cut-coordinate bijection and $U$ the
    unitary map $U\psi=\psi\circ E^{-1}$. For chain vectors $\phi,\psi$ and
    an operator $M$,
    \begin{align}
        \langle U\phi,U\psi\rangle
         & =\langle\phi,\psi\rangle,
        \label{eq:chain_gluing_inner} \\
        \langle U\phi,(UMU^*)U\psi\rangle
         & =\langle\phi,M\psi\rangle.
        \label{eq:chain_gluing_energy}
    \end{align}
    Applying~(\ref{eq:chain_gluing_inner}) and~(\ref{eq:chain_gluing_energy})
    to the ground state and the eigenvector
    extension carries the energy and overlap bounds back to the chain.
\end{proof}

\section{Parameters for a polynomial filter}

\begin{theorem}[A joint count and contraction estimate]
    \label{thm:chain_polynomial_filter_parameters}
    \lean{Real.exists_dg_scalar_parameters}
    \leanok
    For every $C_H>0$ and $M\in\R$, there is a real constant
    $A\ge\max\{1,21,M,\log400000\}$ with the following properties.
    For every integer $d\ge2$ and $0<g\le1$, put
    \begin{align}
        u    & =1+\log d+\log(1/g),                   &
        \ell & =\left\lceil Au^2/g\right\rceil,\notag                                           \\
        c_A  & =6+2\log(A+1),                         & \mathcal D & =(d\ell^2)^{22\ell}.\notag
    \end{align}
    Then $Au^2/g\le\ell\le(A+1)u^2/g$, $A\le\ell$,
    $20\ell+2\lfloor\ell^2/\ell\rfloor=22\ell$, and
    \begin{align}
        \log\mathcal D\le22(A+1)c_Ag^{-1}u^3.\notag
    \end{align}
    For every $0<R\le2C_H\ell$, define
    $\Delta=4\exp(-4\ell^2\sqrt{(4g/5)/R})$. One has
    $\mathcal D^{12}\Delta\le1/100000$.
    The same $A$ has the following simultaneous tail property. For all
    $C,c,\varepsilon>0$ and $p\ge0$ satisfying
    $\max\{0,\log(C/\varepsilon)\}+p+c\le cM$, one has
    $C\exp(-c(\ell-g))\le\varepsilon g^p$.
    This is a supporting scalar estimate for the parameter choice
    in~\cite[lines 1010--1024]{Arad2016RigorousRG}.
\end{theorem}
\begin{proof}\leanok
    Since $\log(A+1)=o(\sqrt A)$, one may choose $A$ above all the prescribed
    lower bounds so that
    \begin{align}
        264\bigl(6+2\log(A+1)\bigr)+1
         & \le4\sqrt{\frac{2A}{5C_H}}.\notag
    \end{align}
    The ceiling bounds follow from $u\ge1$ and $g\le1$. They imply
    $\log\ell\le\log(A+1)+2\log u+\log(1/g)$. Since $\log u\le u$,
    one obtains $\log(d\ell^2)\le c_Au$ and
    $12\log \mathcal D\le264\ell c_Au$.
    The upper ceiling bound gives the asserted bound on $\log\mathcal D$.
    Separately, since $\ell>0$, integer division gives
    $\lfloor\ell^2/\ell\rfloor=\ell$, proving the stated exponent identity.

    From $g\ell\ge Au^2$ and $R\le2C_H\ell$, comparison of nonnegative
    squares gives
    $u\sqrt{2A/(5C_H)}\le\ell\sqrt{(4g/5)/R}$. Consequently
    \begin{align}
        12\log \mathcal D-4\ell^2\sqrt{\frac{4g/5}{R}}
                              & \le-\ell u\le-A,\notag                            \\
        \mathcal D^{12}\Delta & \le4e^{-A}\le\frac4{400000}=\frac1{100000}.\notag
    \end{align}
    Finally put $B=\max\{0,\log(C/\varepsilon)\}$. The threshold condition,
    $\ell\ge Au$, and $u\ge1+\log(1/g)$ give
    \begin{align}
        c(\ell-g) & \ge cAu-c\ge(B+p+c)u-c\notag                             \\
                  & \ge B+p\log(1/g)\ge\log(C/\varepsilon)+p\log(1/g).\notag
    \end{align}
    Exponentiation proves the tail estimate with the same $A$.
\end{proof}

\begin{theorem}[One parameter choice for contraction and the physical tail]
    \label{thm:physical_dg_scalar_parameters}
    \lean{Real.exists_physical_dg_scalar_parameters}
    Let
    \begin{align}
        C_* & =288e^{11/8},\notag                              \\
        M_* & =8\bigl(\max\{0,\log(C_*/10^{-6})\}+1+1/8\bigr).
    \end{align}
    There exists a real $A\ge\max\{1,21,M_*,\log400000\}$ such that,
    for every integer $d\ge2$ and $0<g\le1$, the quantities
    \begin{align}
        u          & =1+\log d+\log(1/g), & \ell & =\lceil Au^2/g\rceil,\notag \\
        \mathcal D & =(d\ell^2)^{22\ell}, & c_A  & =6+2\log(A+1)
    \end{align}
    satisfy $21\le\ell$, $2\le\mathcal D$, $A\le\ell$,
    $Au^2/g\le\ell\le(A+1)u^2/g$, and
    \begin{align}
        \log\mathcal D     & \le22(A+1)c_Ag^{-1}u^3,\notag \\
        C_*e^{-(\ell-g)/8} & \le10^{-6}g.
    \end{align}
    For every $0<R\le8\ell$,
    \begin{align}
        \mathcal D^{12}\,4e^{-4\ell^2\sqrt{(4g/5)/R}}\le10^{-5}.
    \end{align}
    The same $A$ gives both estimates for all $d,g,R$.
\end{theorem}
\begin{proof}
    \uses{thm:chain_polynomial_filter_parameters}
    Apply the joint scalar parameter theorem with boundary constant $4$
    and lower threshold $M_*$. Its simultaneous tail estimate applies with
    $C=C_*$, $c=1/8$, error $10^{-6}$ and power $1$ by the definition of
    $M_*$. The ceiling lower bound gives $\ell\ge21$ and $\mathcal D\ge2$.
    This is the scalar choice used for the physical specialization of
    \cite[Theorem t:ais, lines 1010--1024]{Arad2016RigorousRG}.
\end{proof}

\begin{theorem}[Endpoints for the five-region Chebyshev interval]
    \label{thm:physical_chebyshev_endpoints}
    \lean{Real.physical_chebyshev_endpoints}
    For $\ell\ge21$, $0\le E_0\le4$ and $0<g\le1$, put
    $\eta_0=E_0+g/10$, $\eta_1=E_0+9g/10$ and $R=7\ell-\eta_0$.
    Then $\eta_0<\eta_1\le7\ell$, $0<R\le8\ell$,
    $\eta_1-\eta_0=4g/5$, and
    \begin{align}
        \bigl(2e^{-2\ell^2\sqrt{(\eta_1-\eta_0)/R}}\bigr)^2
        =4e^{-4\ell^2\sqrt{(4g/5)/R}}.
    \end{align}
\end{theorem}
\begin{proof}
    Substitute the definitions. The bounds on $E_0,g,\ell$ make the interval
    nonempty and give the upper bound for $R$. Square the exponential.
\end{proof}

\begin{theorem}[Absorbing a collar at square polynomial degree]
    \label{thm:physical_dg_collar_budget}
    \lean{TNLean.MPS.AreaLaw.collar_card_le_square_at_square_degree}
    Let $q\ge2$, $\ell\ge1$, $d\le2\ell$ and $n$ be natural numbers.
    If $n\le((q\ell^2)^{18\ell+2\lfloor\ell^2/\ell\rfloor})^2$, then
    \begin{align}
        nq^{2d}\le((q\ell^2)^{22\ell})^2.
    \end{align}
\end{theorem}
\begin{proof}
    Since $\ell>0$, the original exponent is $20\ell$.
    The collar contributes $q^{2d}\le(q\ell^2)^{4\ell}$, which is absorbed
    by the difference between $40\ell$ and $44\ell$ in the squared bound.
\end{proof}

\begin{theorem}[One scalar choice for all accuracy parameters]
    \label{thm:dg_accuracy_scalar_parameters}
    \lean{Real.exists_dg_accuracy_scalar_parameters}
    Put $C_*=288e^{11/8}$, $C_+=C_*e^{1/8}$ and
    $M_+=8(\max\{0,\log(C_+/10^{-6})\}+1+1/8)$.
    There is a real constant $A\ge\max\{1,21,M_+,\log400000\}$ such that,
    for every integer $m\ge2$ and $0<g\le1$, the quantities
    \begin{align}
        u_m          & =1+\log m+\log(1/g),     & \ell_m & =\lceil Au_m^2/g\rceil,\notag \\
        \mathcal D_m & =(m\ell_m^2)^{22\ell_m}, & c_A    & =6+2\log(A+1)
    \end{align}
    satisfy $21\le\ell_m$, $m\le\mathcal D_m$, $2\le\mathcal D_m$,
    $A\le\ell_m$, $Au_m^2/g\le\ell_m\le(A+1)u_m^2/g$, and
    \begin{align}
        \log\mathcal D_m     & \le22(A+1)c_Ag^{-1}u_m^3,\notag \\
        C_*e^{-(\ell_m-g)/8} & \le\frac{g}{10^6m}.
    \end{align}
    For every $0<R\le8\ell_m$,
    $\mathcal D_m^{12}\,4e^{-4\ell_m^2\sqrt{(4g/5)/R}}\le10^{-5}$.
    One also has $20\ell_m+2\lfloor\ell_m^2/\ell_m\rfloor=22\ell_m$.
    The same $A$ works for all $m,g,R$.
\end{theorem}
\begin{proof}
    \uses{thm:chain_polynomial_filter_parameters}
    Apply the joint scalar theorem with boundary constant $4$ and threshold
    $M_+$. To obtain the accuracy factor, put $g'=g/m$. Then
    $1+\log(1/g')=u_m$, $Au_m\le\ell_m$ and $0<g'\le1$.
    The scalar tail estimate with coefficient $C_+$ gives
    $C_+e^{-(\ell_m-g')/8}\le10^{-6}g'$.
    Since $g-g'\le1$, the desired tail with coefficient $C_*$ is no larger.
    The ceiling and logarithmic estimates are those of the joint scalar
    choice. The separate integer identity follows from $\ell_m>0$.
    This is a numerical refinement for the physical interval families;
    it does not change the on-site Hilbert space.
\end{proof}

\begin{theorem}[A collar count with a larger numerical base]
    \label{thm:collar_larger_count_base}
    \lean{TNLean.MPS.AreaLaw.collar_card_le_square_of_larger_count_base}
    Let $2\le q\le m$, $\ell\ge1$, $d\le2\ell$ and $n$ be natural numbers.
    If $n\le((q\ell^2)^{18\ell+2\lfloor\ell^2/\ell\rfloor})^2$, then
    \begin{align}
        nq^{2d}\le((m\ell^2)^{22\ell})^2.
    \end{align}
    Only the numerical bound uses $m$; the physical collar still has $q$
    states per site.
\end{theorem}
\begin{proof}
    \uses{thm:physical_dg_collar_budget}
    First absorb the collar using the bound with base $q$.
    Then increase this numerical base from $q$ to $m$ in its nonnegative
    integer power.
\end{proof}

\section{Coefficients of ordered affine products}

The following coefficients separate the two tensor factors in the ordered
word argument of
\cite[Proposition~5.7]{Arad2016RigorousRG} and
\cite[Lemma~4.2]{Arad2013AreaLaw}. Their definition does not
involve the matrices on the second factor.

\begin{definition}[Retained word coefficients]
    \label{def:core_word_coefficients}
    \lean{Matrix.coreWordCoeff}
    \leanok
    Let $m$ be a finite set, let $I$ be an arbitrary set, and choose
    $R_a\in\C^{m\times m}$ for $a\in I$. For a finite list $s$ of pairs
    $(K,c)\in\C^{m\times m}\times\C$ and a word $w$ in $I$, define
    $C_R(s,w)\in\C^{m\times m}$ recursively by
    \begin{align}
        C_R(\varnothing,\varnothing) & =\Id_m,\notag                \\
        C_R(\varnothing,aw)          & =0,\notag                    \\
        C_R((K,c)::s,\varnothing)    & =K C_R(s,\varnothing),\notag \\
        C_R((K,c)::s,aw)             & =K C_R(s,aw)+cR_a C_R(s,w).
        \label{eq:core_word_recursion}
    \end{align}
    Here $aw$ is the word obtained by placing $a$ before $w$, and $(K,c)::s$
    places the displayed pair before the list $s$.
\end{definition}

\begin{lemma}[Expansion over interior words]
    \label{thm:core_word_product}
    \lean{Matrix.prod_coreStep_eq_sum}
    \leanok
    \uses{def:core_word_coefficients}
    Let $I,m,n$ be finite sets. Choose
    $R_a\in\C^{m\times m}$ and $N_a\in\C^{n\times n}$ for $a\in I$.
    For a list $s=((K_0,c_0),\ldots,(K_{h-1},c_{h-1}))$,
    \begin{align}
        \prod_{j=0}^{h-1}\left(K_j\otimes \Id_n+
        \sum_{a\in I}(c_jR_a)\otimes N_a\right)
        =\sum_{t=0}^{h}\ \sum_{w\in I^t}
        C_R(s,w)\otimes\left(\prod_{j=0}^{t-1}N_{w_j}\right).
        \label{eq:core_word_product}
    \end{align}
    Every displayed product is taken in increasing order. Empty products are
    identity matrices, and the empty alphabet and empty matrix index sets are
    included. No commutativity assumption is made on the matrices.
\end{lemma}
\begin{proof}\leanok
    \uses{def:core_word_coefficients}
    Introduce a central formal variable $X$ and the matrix polynomial
    \begin{align}
        P_s(X)=\prod_{(K,c)\in s}
        \left(K\otimes \Id_n+X\sum_{a\in I}(cR_a)\otimes N_a\right),
        \label{eq:core_word_polynomial}
    \end{align}
    with the product in list order. Induction on $s$, using the coefficient
    recurrence for left multiplication by a linear polynomial, gives
    \begin{align}
        [X^t]P_s(X)=\sum_{w\in I^t}C_R(s,w)\otimes
        \left(\prod_{j=0}^{t-1}N_{w_j}\right).
        \label{eq:core_word_coefficient}
    \end{align}
    At a positive degree, split each word into its first letter and remaining
    word. The constant term of the new factor contributes
    $K C_R(s,aw)$, while its linear term contributes $cR_a C_R(s,w)$.
    These are the two terms in (\ref{eq:core_word_recursion}).

    The same recursion shows that $C_R(s,w)=0$ when $|w|>|s|$. Thus
    $\deg P_s\le h$. Evaluation at $X=1$ preserves the ordered product because
    the identity matrix commutes with every coefficient. Summing
    (\ref{eq:core_word_coefficient}) over $0\le t\le h$ gives
    (\ref{eq:core_word_product}).
\end{proof}

\section{Support of three-region operator factors}

Let $V$ be a finite set of sites, with local space $\C^q$ at each site.
For $S\subseteq V$, write $\mathcal H_S=(\C^q)^{\otimes S}$.
If $T,E\subseteq V$ are disjoint and $R=V\setminus(T\cup E)$,
restriction of configurations identifies the full space with
$\mathcal H_T\otimes\mathcal H_E\otimes\mathcal H_R$.
Write $P^g$ for an operator in these coordinates. Support on a set $D$
means that the operator acts trivially on the sites outside $D$.
The statements below give supporting locality algebra for
\cite[Proposition~5.7]{Arad2016RigorousRG}.

\begin{theorem}[Support of partial matrix elements]
    \label{thm:three_region_partial_support}
    \lean{TNLean.MPS.AreaLaw.operatorBlock_isSupportedOn_of_threeSplit_left,
        TNLean.MPS.AreaLaw.leftBlock_isSupportedOn_of_threeSplit_right}
    \leanok
    Suppose that $P$ is supported on $D\subseteq V$.
    If $P^g=U\otimes \Id_R$, then for any configurations $a,b$ on $T$,
    the matrix on $E$ with entries
    \begin{align}
        U^{\mathrm R}_{ab}(x,y)=U((a,x),(b,y))
        \label{eq:three_region_left_partial}
    \end{align}
    is supported on $D\cap E$, viewed as a subset of $E$.
    If instead $P^g=\Id_T\otimes V_0$, then for any configurations $r,s$
    on $R$, the matrix
    \begin{align}
        (V_0)^{\mathrm L}_{rs}(x,y)=V_0((x,r),(y,s))
        \label{eq:three_region_right_partial}
    \end{align}
    has the same support conclusion. The two fixed configurations in either
    statement need not agree outside $D$.
\end{theorem}
\begin{proof}\leanok
    For $q>0$, choose one fixed configuration on the unused region.
    The displayed partial matrix element is then an entry of $P$ with that
    configuration on both sides. A disagreement at a site of $E\setminus D$
    forces the entry to vanish. When the configurations agree there, the
    dependence on the remaining coordinates follows from the support of $P$.
    For $q=0$, the configuration set of $E$ is either empty or a singleton;
    the support assertion follows directly in both cases.
\end{proof}

\begin{theorem}[Equality from a middle identity lift]
    \label{thm:three_region_middle_lift_equality}
    \lean{TNLean.MPS.AreaLaw.eq_of_threeSplit_middle_lift_eq}
    \leanok
    Suppose that $E$ contains a specified site. If $W,M$ are operators on
    $\mathcal H_E$ and
    \begin{align}
        \Id_T\otimes W\otimes \Id_R=\Id_T\otimes M\otimes \Id_R,
        \label{eq:three_region_middle_lift_equality}
    \end{align}
    then $W=M$. No positive-local-dimension assumption is required.
\end{theorem}
\begin{proof}\leanok
    For $q>0$, evaluate the identity at equal fixed configurations on $T$ and
    $R$. For $q=0$, a configuration on $E$ would assign an element of the empty
    local alphabet to the specified site. Thus both matrices have empty index
    sets and are equal.
\end{proof}

\begin{theorem}[Canonical placement in the middle region]
    \label{thm:three_region_middle_embedding}
    \lean{TNLean.MPS.AreaLaw.embedOp_subtype_val_submatrix_threeSplit}
    \leanok
    For every operator $W$ on $\mathcal H_E$, its canonical extension to $V$
    has three-region form
    \begin{align}
        (\iota_E(W))^g=\Id_T\otimes W\otimes \Id_R.
        \label{eq:three_region_middle_embedding}
    \end{align}
    This includes empty regions and zero local dimension.
\end{theorem}
\begin{proof}\leanok
    The extension has entry $W(x,y)$ when the two configurations agree off
    $E$, and zero otherwise. In three-region coordinates, agreement off $E$
    is precisely equality of both the $T$ and $R$ configurations. These two
    equalities give the entries of the two identity matrices.
\end{proof}

\begin{theorem}[Two-window factors with their collars]
    \label{thm:three_region_two_window_collars}
    \lean{TNLean.MPS.AreaLaw.exists_threeSplit_two_window_factors_with_collars}
    \leanok
    \uses{def:chain_model,lem:peps_frame_split_forms,
        thm:three_region_partial_support}
    Consider a chain of $L$ sites and nonnegative integers $a,b,\ell$ with
    $a+2\ell\le b$ and $b+2\ell\le L-1$.
    Let the label set consist of $L_r,R_r$ for $0\le r<2\ell$, together
    with $B_0,B_1,B_2$. Suppose the supplied operator family $F$ satisfies
    \begin{align}
        \operatorname{supp}F_{L_r} & \subseteq[a+r,a+r+2),         &
        \operatorname{supp}F_{R_r} & \subseteq[b+r,b+r+2),\notag     \\
        \operatorname{supp}F_{B_0} & \subseteq[0,a+1),             &
        \operatorname{supp}F_{B_1} & \subseteq[a+2\ell,b+1),\notag   \\
        \operatorname{supp}F_{B_2} & \subseteq[b+2\ell,L).
        \label{eq:three_region_window_supports}
    \end{align}
    Intervals are intersected with the finite chain. Fix $0\le p,s<2\ell$
    and set
    \begin{align}
        T & =[0,a+p+1),\notag           \\
        E & =[a+p+1,b+s+1),\notag       \\
        R & =V\setminus(T\cup E),\notag \\
        C & =E\setminus[a+2\ell+1,b).
        \label{eq:three_region_selected_regions}
    \end{align}
    There are regional families $H^L,H^M,H^R$ and crossing operators $U,V_0$
    such that, for every label $i\ne L_p,R_s$,
    \begin{align}
        F_i^g                   & =H_i^L\otimes \Id_E\otimes \Id_R
        +\Id_T\otimes H_i^M\otimes \Id_R
        +\Id_T\otimes \Id_E\otimes H_i^R,\notag                      \\
        F_{L_p}^g               & =U\otimes \Id_R,                 &
        F_{R_s}^g               & =\Id_T\otimes V_0,\notag           \\
        \operatorname{OSR}(U)   & \le q^2,                         &
        \operatorname{OSR}(V_0) & \le q^2.
        \label{eq:three_region_two_window_factors}
    \end{align}
    The same choices also satisfy
    \begin{align}
        H^M_{B_0}                    & =H^M_{B_2}=0,                   &
        H^L_{B_1}                    & =H^R_{B_1}=0,\notag               \\
        H^M_{L_r}                    & =0\quad(r\le p),                &
        H^M_{R_r}                    & =0\quad(s\le r),\notag            \\
        \operatorname{supp}H^M_{L_r} & \subseteq C,                    &
        \operatorname{supp}H^M_{R_r} & \subseteq C\quad(0\le r<2\ell).
        \label{eq:three_region_collar_conditions}
    \end{align}
    The three bulk operators are supplied independently; no identity with
    complementary sums of bond terms is assumed.
\end{theorem}
\begin{proof}\leanok
    \uses{thm:chain_two_cut_edge_supports,lem:peps_frame_split_forms,
        thm:peps_supported_left_factor_rank,thm:peps_supported_right_factor_rank,
        thm:three_region_partial_support}
    The selected bonds meet each side of their cut at one site, so their
    crossing operators have operator Schmidt rank at most $q^2$.
    Every other left-window bond lies either in $T$ or in $E$, according as
    its position is before or after $p$. Similarly, every other right-window
    bond lies in $E$ or $R$ according as its position is before or after $s$.
    Choose zero for every unused regional component, including the two selected
    middle components. The three bulk supports lie in $T,E,R$, respectively.
    This gives the identities and zero conditions.

    Each nonzero window middle component inherits its original two-site support
    restricted to $E$, by the partial-matrix-element argument above applied to
    an identity lift. The bounds on the window positions place this restricted
    support in $C$. Monotonicity of support gives the last two assertions.
\end{proof}

\begin{theorem}[The weighted middle sum and its supported remainder]
    \label{thm:three_region_weighted_middle_remainder}
    \lean{TNLean.MPS.AreaLaw.exists_weighted_ordinary_middle_eq_central_add_supported}
    \leanok
    Let $E$ be any finite set, let $H^M$ be indexed by the labels above, and
    let $C\subseteq E$. Suppose $H^M_{B_0}=H^M_{B_2}=0$,
    $H^M_{L_r}=0$ for $r\le p$, and $H^M_{R_r}=0$ for $s\le r$.
    Suppose also that $H^M_{L_r}$ is supported on $C$ for $p<r$, and that
    $H^M_{R_r}$ is supported on $C$ for $r<s$.
    For arbitrary complex weights $z_i$, put
    $O=\{i:i\ne L_p,R_s\}$. There is an operator $K$ supported on $C$ with
    \begin{align}
        \sum_{i\in O}z_iH_i^M=z_{B_1}H^M_{B_1}+K.
        \label{eq:three_region_weighted_middle_remainder}
    \end{align}
\end{theorem}
\begin{proof}\leanok
    Take
    \begin{align}
        K=\sum_{p<r<2\ell}z_{L_r}H^M_{L_r}
        +\sum_{0\le r<s}z_{R_r}H^M_{R_r}.
        \label{eq:three_region_weighted_remainder_formula}
    \end{align}
    The discarded terms vanish by the stated hypotheses. Support on $C$ is
    preserved by scalar multiplication and finite addition.
\end{proof}

\section{Prescribed middle factors of polynomial coefficients}

The algebraic statements in this section describe the ordered
products used in \cite[Proposition~5.7]{Arad2016RigorousRG}.
The middle operators are specified before the coefficients on the exterior
factor. All products retain their order.
An interval $\{0,\ldots,h-1\}$ is empty when $h=0$.

For a matrix $M$ indexed by a product $m\times n$, use the right operator
blocks and left coefficient slices
\begin{align}
    M^{\mathrm R}_{ab}(x,y) & =M((a,x),(b,y)),\notag \\
    M^{\mathrm L}_{xy}(a,b) & =M((a,x),(b,y)).
    \label{eq:two_marker_operator_blocks}
\end{align}
Here $a,b\in m$ and $x,y\in n$; the right blocks are indexed by $n$
and the left slices by $m$.

\begin{lemma}[Prescribed factors from operator blocks]
    \label{thm:two_marker_prescribed_span}
    \lean{Matrix.exists_left_factors_of_operatorBlock_mem_span,Matrix.exists_right_factors_of_left_blocks_mem_span}
    \leanok
    Let $\mathcal I$ be finite, let $m,n$ be arbitrary sets, and let
    $C_\alpha$ be complex matrices indexed by $n$. Suppose that
    $M^{\mathrm R}_{ab}\in\spn_{\C}\{C_\alpha:\alpha\in\mathcal I\}$ for every $a,b\in m$.
    There are matrices $L_\alpha$ indexed by $m$ such that
    \begin{align}
        M=\sum_{\alpha\in\mathcal I}L_\alpha\otimes C_\alpha.
        \label{eq:two_marker_prescribed_span}
    \end{align}
    For the opposite orientation, let $C_\alpha\in\C^{m\times m}$ for
    $\alpha\in\mathcal I$. If
    $M^{\mathrm L}_{xy}\in\spn_{\C}\{C_\alpha:\alpha\in\mathcal I\}$
    for every $x,y\in n$, then
    $M=\sum_{\alpha\in\mathcal I}C_\alpha\otimes R_\alpha$ for suitable
    matrices $R_\alpha$ indexed by $n$.
\end{lemma}
\begin{proof}\leanok
    For each $a,b$, choose coefficients $\lambda_{ab,\alpha}$ with
    $M^{\mathrm R}_{ab}=\sum_\alpha\lambda_{ab,\alpha}C_\alpha$, and set
    $(L_\alpha)_{ab}=\lambda_{ab,\alpha}$. Comparing matrix entries gives
    (\ref{eq:two_marker_prescribed_span}). The opposite orientation follows by
    the same argument with the two pairs of indices exchanged.
\end{proof}

\begin{definition}[Crossing words and ordinary run lengths]
    \label{def:two_marker_middle_index}
    \lean{TNLean.MPS.AreaLaw.TwoMarkerMiddleIndex}
    \leanok
    For nonnegative integers $r,i,j,u,v$, let $\mathcal I_{r,i,j}^{u,v}$ consist
    of pairs $(w,g)$ with
    \begin{align}
        w & :\{0,\ldots,i+j-1\}\longrightarrow
        \{0,\ldots,u-1\}\sqcup\{0,\ldots,v-1\},\notag       \\
        g & :\{0,\ldots,i+j\}\longrightarrow\{0,\ldots,r\}.
        \label{eq:two_marker_index_data}
    \end{align}
    The word has $i$ letters in the first summand and $j$ in the second, with
    $\sum_{h=0}^{i+j}g_h+i+j\le r$.
\end{definition}

\begin{definition}[Ordered middle product]
    \label{def:two_marker_middle_product}
    \lean{TNLean.MPS.AreaLaw.twoMarkerMiddleFactor}
    \leanok
    \uses{def:two_marker_middle_index}
    Let $A$ be an element of a monoid, and let $U_0,\ldots,U_{u-1}$ and
    $V_0,\ldots,V_{v-1}$ belong to the same monoid. For
    $(w,g)\in\mathcal I_{r,i,j}^{u,v}$, put $T(\operatorname{inl}s)=U_s$,
    $T(\operatorname{inr}t)=V_t$, and define
    \begin{align}
        \Phi_{w,g}(A,U,V)
        =\left(\prod_{h=0}^{i+j-1} A^{g_h}T(w_h)\right)A^{g_{i+j}}.
        \label{eq:two_marker_middle_product}
    \end{align}
    The product over $h$ is in increasing order; an empty product is the
    identity.
\end{definition}

\begin{definition}[Two distinguished frequencies]
    \label{def:two_marker_coefficient}
    \lean{TNLean.MPS.AreaLaw.twoMarkerCoefficient}
    \leanok
    Let $D,U,V$ belong to a semiring. With two central formal variables $x,y$,
    define
    \begin{align}
        C_{r,i,j}(D,U,V)=[x^i y^j](D+yV+xU)^r.
        \label{eq:two_marker_coefficient}
    \end{align}
    The variable $x$ counts $U$ and $y$ counts $V$; the coefficients $D,U,V$
    need not commute.
\end{definition}

\begin{lemma}[Operator blocks of sums and tensor products]
    \label{thm:two_marker_block_identities}
    \lean{TNLean.MPS.AreaLaw.operatorBlock_sum,TNLean.MPS.AreaLaw.operatorBlock_kronecker_eq_smul}
    \leanok
    For a finite family of matrices $M_\alpha$ indexed by $m\times n$, and
    matrices $E$ and $W$ indexed by $m$ and $n$, respectively,
    \begin{align}
        \left(\sum_\alpha M_\alpha\right)^{\mathrm R}_{ab}
                                      & =\sum_\alpha(M_\alpha)^{\mathrm R}_{ab},\notag \\
        (E\otimes W)^{\mathrm R}_{ab} & =E_{ab}W.
        \label{eq:two_marker_block_identities}
    \end{align}
    No finiteness assumption on $m$ or $n$ is required.
\end{lemma}
\begin{proof}\leanok
    Both identities follow by comparing the entry indexed by $x,y\in n$.
\end{proof}

\begin{lemma}[A common prescribed middle family through a degree bound]
    \label{thm:two_marker_common_middle}
    \lean{TNLean.MPS.AreaLaw.exists_sum_twoMarkerCoefficient_decomposition}
    \leanok
    \uses{def:two_marker_middle_index,def:two_marker_middle_product,def:two_marker_coefficient}
    Let $m,n$ be finite sets. Choose complex matrices $E,E_s^U,E_t^V$ indexed
    by $m$ and $A,U_s^M,V_t^M$ indexed by $n$, with $0\le s<u$ and $0\le t<v$.
    Put
    \begin{align}
        D & =E\otimes I+I\otimes A,\notag              \\
        U & =\sum_{s=0}^{u-1}E_s^U\otimes U_s^M,\notag \\
        V & =\sum_{t=0}^{v-1}E_t^V\otimes V_t^M.
        \label{eq:two_marker_split_letters}
    \end{align}
    For all nonnegative integers $k,i,j$ and all complex numbers
    $c_0,\ldots,c_k$, there are matrices $B_{w,g}$ indexed by $m$ such that
    \begin{align}
        \sum_{r=0}^k c_r C_{r,i,j}(D,U,V)
        =\sum_{(w,g)\in\mathcal I_{k,i,j}^{u,v}}
        B_{w,g}\otimes\Phi_{w,g}(A,U^M,V^M).
        \label{eq:two_marker_common_middle}
    \end{align}
    The displayed middle family depends only on $A,U^M,V^M,k,i,j$.
\end{lemma}
\begin{proof}\leanok
    \uses{thm:two_marker_prescribed_span,thm:two_marker_block_identities,def:two_marker_middle_index,def:two_marker_middle_product,def:two_marker_coefficient}
    Induction on $r$ shows that every right operator block of
    $C_{r,i,j}(D,U,V)$ belongs to the span of the products indexed by
    $\mathcal I_{r,i,j}^{u,v}$. At degree zero, the coefficient is $I$ for
    $i=j=0$ and zero otherwise. The positive-frequency recurrence is
    \begin{align}
        C_{r+1,i+1,j+1}
        =C_{r,i+1,j+1}D+C_{r,i+1,j}V+C_{r,i,j+1}U.
        \label{eq:two_marker_right_recurrence}
    \end{align}
    When one frequency vanishes, the corresponding last-letter term is absent.
    Right multiplication by $E\otimes I$ preserves a middle word and enlarges
    its degree bound. Right multiplication by $I\otimes A$ increases its final
    ordinary run. The $U$ and $V$ terms append the corresponding crossing
    letter, retaining the order in (\ref{eq:two_marker_middle_product}).
    Lemma~\ref{thm:two_marker_prescribed_span} and
    (\ref{eq:two_marker_block_identities}) justify these span operations.

    If $r\le k$, the same crossing letters and run lengths define an element
    of $\mathcal I_{k,i,j}^{u,v}$ with the same ordered product. Thus all the
    right operator blocks in the left side of
    (\ref{eq:two_marker_common_middle}) belong to the span of the one displayed
    middle family. Apply Lemma~\ref{thm:two_marker_prescribed_span} once more
    to obtain the matrices $B_{w,g}$.
\end{proof}

\begin{lemma}[Number of prescribed middle factors]
    \label{thm:two_marker_middle_cardinality}
    \lean{TNLean.MPS.AreaLaw.card_twoMarkerMiddleIndex_square_le}
    \leanok
    \uses{def:two_marker_middle_index}
    For all nonnegative integers $r,i,j,q$,
    \begin{align}
        |\mathcal I_{r,i,j}^{q^2,q^2}|
        \le 2^{i+j}q^{2(i+j)}(r+1)^{i+j+1}.
        \label{eq:two_marker_middle_cardinality}
    \end{align}
\end{lemma}
\begin{proof}\leanok
    \uses{def:two_marker_middle_index}
    There are at most $(2q^2)^{i+j}$ choices for the decorated word and
    $(r+1)^{i+j+1}$ choices for its run lengths. Discarding the frequency and
    total-length restrictions gives their product as an upper bound.
\end{proof}

\section{A two-binomial count for crossing words}

\begin{theorem}[Counting crossing decorations and ordinary runs]
    \label{thm:two_marker_index_two_binomial}
    \lean{TNLean.MPS.AreaLaw.card_twoMarkerMiddleIndex_two_binomial_le}
    \uses{def:two_marker_middle_index}
    Let $k,i,j,u,v$ be nonnegative integers, and let
    $\mathcal I_{k,i,j}^{u,v}$ be the set in
    Definition~\ref{def:two_marker_middle_index}. Then
    \begin{align}
        \bigl|\mathcal I_{k,i,j}^{u,v}\bigr|
        \le u^i v^j
        \binom{k+i+1}{2i+1}\binom{k+j+1}{2j+1}.
    \end{align}
    In particular, no positivity assumption on $k,u,v$ is required.
    This is the finite count used in
    \cite[Proposition~5.7]{Arad2016RigorousRG}.
\end{theorem}
\begin{proof}
    Put $m=i+j$. If $m>k$, the set is empty. Suppose $m\le k$.
    The left positions form an $i$-element subset of $\{0,\ldots,m-1\}$.
    Recording this subset and the labels in its two complementary parts
    determines the crossing word. There are at most $\binom{m}{i}u^i v^j$
    choices. Adjoining the unused degree to the $m+1$ ordinary run lengths
    gives a weak composition of $k-m$ into $m+2$ parts. Deleting the added
    coordinate recovers the run lengths, so their number is at most
    $\binom{k+1}{m+1}$. Therefore
    \begin{align}
        \bigl|\mathcal I_{k,i,j}^{u,v}\bigr|
        \le u^i v^j\binom{k+1}{m+1}\binom{m}{i}.
    \end{align}
    Pascal's identity and the binomial product identity give
    \begin{align}
        \binom{k+1}{m+1}\binom{m}{i}
         & \le\binom{k+1}{m+1}\binom{m+1}{i+1}
        =\binom{k+1}{i+1}\binom{k-i}{j} \notag  \\
         & \le\binom{k+1}{i+1}\binom{k+1}{j+1}.
    \end{align}
    Symmetry and monotonicity in the upper argument yield
    \begin{align}
        \binom{k+1}{i+1}
        =\binom{k+1}{k-i}
        \le\binom{k+i+1}{k-i}
        =\binom{k+i+1}{2i+1}.
    \end{align}
    Applying the same inequality to $j$ proves the bound. Empty label sets
    are included, with $0^0=1$.
\end{proof}


\begin{theorem}[A square bound for the middle-family cardinality]
    \label{thm:chain_two_window_family_cardinality}
    \lean{TNLean.MPS.AreaLaw.crossing_grid_cardinality_le_square}
    \leanok
    Let $q,k,\ell,Z,N$ be nonnegative integers with $q\ge2$, $k\ge1$ and
    $\ell\ge1$, and put $m=\lfloor k/\ell\rfloor$ and
    $r=\lfloor k/(2\ell)\rfloor$.
    Suppose that $Z\le(k+1)^{4\ell+3}$ and
    \begin{align}
        N & \le (2\ell)^2(r+1)^2 Z(qk)^{24\ell+4m}.\label{eq:chain_middle_family_count}
    \end{align}
    Then $N\le D^2$, where $D=(qk)^{18\ell+2m}$.
    This is a numerical estimate for the counting argument supporting
    \cite[Proposition~5.7]{Arad2016RigorousRG}.
\end{theorem}
\begin{proof}\leanok
    The inequalities $2\ell\le2^\ell\le(qk)^\ell$ and
    $r+1\le k+1\le qk\le(qk)^\ell$ give
    $(2\ell)^2(r+1)^2\le(qk)^{4\ell}$.
    Also $Z\le(qk)^{4\ell+3}\le(qk)^{7\ell}$.
    Substituting into~(\ref{eq:chain_middle_family_count}) yields
    \begin{align*}
        N\le(qk)^{35\ell+4m}\le(qk)^{36\ell+4m}=D^2.
    \end{align*}
\end{proof}

\section{Fixed middle data for selected windows}

Let the chain have sites $\{0,\ldots,L-1\}$ and local space $\C^q$.
All intervals below are intersected with this chain. For nonnegative
integers $a,b,\ell$ and positions $0\le p,s<2\ell$, put
\begin{align}
    K & =[a+2\ell,b+1),                      & Z &
    =[a+2\ell+1,b),\notag                                         \\
    T & =[0,a+p+1),                          & E &
    =[a+p+1,b+s+1),\notag                                         \\
    R & =\{0,\ldots,L-1\}\setminus(T\cup E), & C & =E\setminus Z.
    \label{eq:selected_middle_regions}
\end{align}
These statements provide supporting locality algebra for
\cite[Proposition~5.7]{Arad2016RigorousRG}.

\begin{lemma}[The central regions lie in the selected middle]
    \label{thm:selected_middle_core_inclusions}
    \lean{TNLean.MPS.AreaLaw.core_and_strict_core_subset_selected_middle}
    \leanok
    For the data above, $K\subseteq E$ and $Z\subseteq E$.
    No separation assumption on $a,b$ is needed for these inclusions.
\end{lemma}
\begin{proof}\leanok
    The inequalities $p<2\ell$ and $s\ge0$ show that the left endpoint of
    $E$ is no greater than that of $K$, while its right endpoint is no
    smaller. The same comparison applies to $Z$. Monotonicity of clipped
    intervals proves the assertions.
\end{proof}

\begin{lemma}[Cardinality of the selected collars]
    \label{thm:selected_middle_collar_cardinality}
    \lean{TNLean.MPS.AreaLaw.card_selected_middle_sdiff_strict_core_le}
    \leanok
    If $a+2\ell\le b$, then $|C|\le4\ell$.
\end{lemma}
\begin{proof}\leanok
    The complement of the strict core within $E$ is
    \begin{align}
        C=[a+p+1,a+2\ell+1)\cup
        [b,b+s+1).
        \label{eq:selected_middle_collar_union}
    \end{align}
    The two intervals may overlap when the windows touch. Their cardinalities
    are at most $2\ell-p$ and $s+1$, respectively. Both are at most $2\ell$,
    so the cardinality of their union is at most $4\ell$.
\end{proof}

\begin{lemma}[Middle and crossing factors chosen before the exterior operators]
    \label{thm:selected_middle_fixed_physical_data}
    \lean{TNLean.MPS.AreaLaw.exists_physical_selected_window_fixed_middle_data}
    \leanok
    \uses{thm:selected_middle_core_inclusions,thm:three_region_partial_support,
        thm:three_region_middle_lift_equality,thm:three_region_middle_embedding,
        thm:three_region_two_window_collars,thm:three_region_weighted_middle_remainder,
        thm:two_marker_prescribed_span}
    Assume $a+2\ell\le b$ and $b+2\ell\le L-1$.
    Let $W^L_r,W^R_r$, for $0\le r<2\ell$, be operators on the full chain
    supported on $[a+r,a+r+2)$ and
    $[b+r,b+r+2)$, respectively.
    Let $M$ be an arbitrary operator on $K$. Write $P_M$ for its canonical
    extension to the full chain and $M_E$ for its canonical extension to $E$.

    Use labels $L_r,R_r,B_0,B_1,B_2$, and let $O$ be all labels except
    $L_p,R_s$. There exist sites $e_L,e_R\in E$, a middle family $H^M$,
    crossing operators $U,V_0$, and exterior families
    $A_z\in\operatorname{End}(\mathcal H_T)$,
    $B_z\in\operatorname{End}(\mathcal H_R)$, indexed by
    $z\in\{0,\ldots,q^2-1\}$, with the following properties.

    First, the sites are the literal selected endpoints:
    \begin{align}
        e_L=a+p+1,\qquad e_R=b+s,
        \label{eq:selected_middle_endpoints}
    \end{align}
    and the central middle factor is $H^M_{B_1}=M_E$.
    For $z$ corresponding to a pair $(i,j)\in\{0,\ldots,q-1\}^2$,
    let $N_e(z)$ be the one-site matrix unit $|i\rangle\langle j|$
    canonically extended from $e$ to $E$.
    In three-region coordinates the actual selected crossings satisfy
    \begin{align}
        (W^L_p)^g & =U\otimes I_R,                 &
        (W^R_s)^g & =I_T\otimes V_0,\notag           \\
        U         & =\sum_z A_z\otimes N_{e_L}(z), &
        V_0       & =\sum_z N_{e_R}(z)\otimes B_z.
        \label{eq:selected_middle_fixed_crossings}
    \end{align}
    Every $N_{e_L}(z)$ and $N_{e_R}(z)$ is supported on $C$.
    Moreover, for every complex weight family $z_i$ there is an operator
    $K_z$ supported on $C$ such that
    \begin{align}
        \sum_{i\in O}z_iH^M_i=z_{B_1}M_E+K_z.
        \label{eq:selected_middle_weighted_sum}
    \end{align}

    All the preceding choices are made before the following exterior data.
    For every pair of full-chain operators $P_L,P_R$ supported on
    $[0,a+1)$ and
    $[b+2\ell,L)$, respectively, set
    \begin{align}
        F_{L_r}=W^L_r,\quad F_{R_r}=W^R_r,\quad
        F_{B_0}=P_L,\quad F_{B_1}=P_M,\quad F_{B_2}=P_R.
        \label{eq:selected_middle_full_family}
    \end{align}
    There are exterior regional families $H^L,H^R$ such that, for $i\in O$,
    \begin{align}
        F_i^g=H_i^L\otimes I_E\otimes I_R
        +I_T\otimes H_i^M\otimes I_R
        +I_T\otimes I_E\otimes H_i^R,
        \label{eq:selected_middle_late_exterior}
    \end{align}
    with $H^L_{B_1}=H^R_{B_1}=0$.
    The sites $e_L,e_R$ may coincide. No positive-local-dimension assumption
    is imposed.
\end{lemma}
\begin{proof}\leanok
    \uses{thm:selected_middle_core_inclusions,thm:three_region_partial_support,
        thm:three_region_middle_lift_equality,thm:three_region_middle_embedding,
        thm:three_region_two_window_collars,thm:three_region_weighted_middle_remainder,
        thm:two_marker_prescribed_span}
    Apply the three-region collar decomposition with both exterior bulk
    operators set to zero and with central operator $P_M$.
    The two selected bonds meet $E$ only at $a+p+1$ and $b+s$.
    Both endpoints lie in $C$, as follows directly from $p,s<2\ell$.

    Canonical placement of $M$ factors through its inclusion into $E$.
    Since the central exterior components vanish, the chosen central middle
    factor satisfies
    \begin{align}
        P_M^g=I_T\otimes H^M_{B_1}\otimes I_R
        =I_T\otimes M_E\otimes I_R.
        \label{eq:selected_middle_central_lift}
    \end{align}
    The endpoint $e_L$ supplies an actual site of $E$, so the two middle
    factors in (\ref{eq:selected_middle_central_lift}) coincide,
    including at $q=0$.

    The partial matrix elements of each crossing are supported on its single
    middle endpoint. An operator supported on one site belongs to the span of
    the canonically extended matrix units at that site. Retaining these
    prescribed spanning families in the two product decompositions gives
    (\ref{eq:selected_middle_fixed_crossings}). Their support lies in $C$.
    The weighted supported-remainder identity then gives
    (\ref{eq:selected_middle_weighted_sum}).

    Finally, insert arbitrary $P_L,P_R$ with the stated supports by changing
    only the corresponding exterior bulk components. The middle family and
    the exterior components of the window labels remain fixed. This proves
    (\ref{eq:selected_middle_late_exterior}) with the central exterior
    components equal to zero.
\end{proof}

\section{Prescribed middle products of selected coefficients}

The following decomposition is supporting algebra for
\cite[Proposition~5.7]{Arad2016RigorousRG}.
It retains the actual central operator and the literal one-site matrix
units furnished by the selected crossing edges.

\begin{definition}[Grouping the exterior regions]
    \label{def:selected_exterior_middle_equiv}
    \lean{TNLean.MPS.AreaLaw.exteriorMiddleEquiv}
    \leanok
    For sets $X,Y,Z$, define the equivalence
    \begin{align}
        \epsilon:(X\times Y)\times Z & \longrightarrow(X\times Z)\times Y, &
        \epsilon((x,y),z)            & =((x,z),y).
        \label{eq:selected_exterior_middle_equiv}
    \end{align}
    This changes only the grouping of the three factors; the coordinates
    within each factor retain their original order.
\end{definition}

\begin{theorem}[One prescribed middle family for all exterior operators]
    \label{thm:selected_coefficient_prescribed_middle}
    \lean{TNLean.MPS.AreaLaw.exists_physical_twoMarkerCoefficient_prescribed_middle_decomposition}
    \leanok
    \uses{def:selected_exterior_middle_equiv,thm:selected_middle_core_inclusions,
        thm:selected_middle_fixed_physical_data,def:two_marker_middle_index,
        def:two_marker_middle_product,def:two_marker_coefficient,
        thm:two_marker_common_middle}
    Let $L,q,a,b,\ell$ be nonnegative integers with
    $a+2\ell\le b$ and $b+2\ell\le L-1$.
    Fix positions $0\le p,s<2\ell$ and full-chain operators $W^L_r,W^R_r$
    supported on $[a+r,a+r+2)$ and
    $[b+r,b+r+2)$, respectively.
    Let $M$ be an operator on $K=[a+2\ell,b+1)$.
    Every interval is intersected with the chain of $L$ sites.
    Put
    \begin{align}
        T & =[0,a+p+1),                          & E &
        =[a+p+1,b+s+1),\notag                          \\
        R & =\{0,\ldots,L-1\}\setminus(T\cup E), &
        C & =E\setminus[a+2\ell+1,b).
        \label{eq:selected_coefficient_regions}
    \end{align}
    Let $M_E$ and $P_M$ be the canonical extensions of $M$ to $E$ and to the
    full chain. Use labels $L_r,R_r,B_0,B_1,B_2$, and let $O$ exclude precisely
    $L_p,R_s$. Fix arbitrary complex weights $z_i$ and nonnegative integers
    $k,i,j$.

    There are sites $e_L,e_R\in E$ and middle operators $A,K_z$ such that
    \begin{align}
        e_L & =a+p+1,          & e_R                    & =b+s,\notag  \\
        A   & =z_{B_1}M_E+K_z, & \operatorname{supp}K_z & \subseteq C.
        \label{eq:selected_coefficient_middle_data}
    \end{align}
    For each $e\in E$, let $N_e(t)$, $0\le t<q^2$, be the canonical
    extensions to $E$ of the one-site matrix units, indexed by pairs of local
    labels. Every $N_{e_L}(t)$ and $N_{e_R}(t)$ is supported on $C$.
    Define the prescribed middle family by
    \begin{align}
        \Psi_f=\Phi_f(A,N_{e_L},N_{e_R}),
        \qquad f\in\mathcal I_{k,i,j}^{q^2,q^2},
        \label{eq:selected_coefficient_middle_family}
    \end{align}
    where $\Phi$ is the ordered middle product.

    These sites and operators are chosen before the following data.
    For every pair $P_L,P_R$ supported on $[0,a+1)$ and
    $[b+2\ell,L)$,
    set
    \begin{align}
        F_{L_r}           & =W^L_r,               & F_{R_r} & =W^R_r,\notag                         \\
        F_{B_0}           & =P_L,                 & F_{B_1} & =P_M,         & F_{B_2} & =P_R,\notag \\
        D_{\mathrm{phys}} & =\sum_{t\in O}z_tF_t.
        \label{eq:selected_coefficient_physical_sum}
    \end{align}
    For every complex family $c_0,\ldots,c_k$, there are exterior operators
    $B_f$ on $\mathcal H_T\otimes\mathcal H_R$ for which
    \begin{align}
        \left(\sum_{r=0}^k c_r
        C_{r,i,j}(D_{\mathrm{phys}},W^L_p,W^R_s)\right)^{g;\epsilon}
        =\sum_{f\in\mathcal I_{k,i,j}^{q^2,q^2}}B_f\otimes\Psi_f.
        \label{eq:selected_coefficient_decomposition}
    \end{align}
    Here $g$ restricts a full configuration to the three regions, after which
    $\epsilon$ groups the two exterior configurations together. The notation
    $C_{r,i,j}$ denotes the two distinguished-frequency coefficient. The
    selected crossings $W^L_p,W^R_s$ are unweighted; the weights occur only
    in $D_{\mathrm{phys}}$.
    The assertion includes zero local dimension and every value of $k,i,j$.
\end{theorem}
\begin{proof}\leanok
    \uses{def:selected_exterior_middle_equiv,
        thm:selected_middle_fixed_physical_data,thm:two_marker_common_middle}
    Choose the fixed middle and crossing data from the two windows and $M$.
    Their ordinary weighted middle sum is $A=z_{B_1}M_E+K_z$, with $K_z$
    supported on $C$. Retain the prescribed one-site matrix-unit families in
    both crossing decompositions.

    After $P_L,P_R$ are supplied, sum their regional identities with the
    weights $z_t$. In three-region coordinates the result has ordinary part
    \begin{align}
        D_0=L_z\otimes I_E\otimes I_R
        +I_T\otimes A\otimes I_R
        +I_T\otimes I_E\otimes R_z.
        \label{eq:selected_coefficient_ordinary_form}
    \end{align}
    The crossings have the forms in
    (\ref{eq:selected_middle_fixed_crossings}), with exterior factors
    $A^L_t$ on $T$ and $B^R_t$ on $R$:
    $U=\sum_tA^L_t\otimes N_{e_L}(t)$ and
    $V_0=\sum_tN_{e_R}(t)\otimes B^R_t$. Regrouping by $\epsilon$ gives
    \begin{align}
        (D_0)^\epsilon            & =(L_z\otimes I_R+I_T\otimes R_z)\otimes I_E
        +I_{T\times R}\otimes A,\notag                                           \\
        (U\otimes I_R)^\epsilon   & =\sum_t(A^L_t\otimes I_R)\otimes N_{e_L}(t),
        \notag                                                                   \\
        (I_T\otimes V_0)^\epsilon & =\sum_t(I_T\otimes B^R_t)\otimes N_{e_R}(t).
        \label{eq:selected_coefficient_regrouped_form}
    \end{align}
    Thus the exterior factors are combined, while the middle factors $A$,
    $N_{e_L}(t)$ and $N_{e_R}(t)$ are unchanged.

    Reindexing is a ring homomorphism, so it commutes with powers, extraction
    of the two marked coefficients, and their finite scalar combinations.
    Apply the prescribed-middle decomposition through degree $k$ to these
    regrouped operators in (\ref{eq:selected_coefficient_regrouped_form}). Its middle factors are exactly
    (\ref{eq:selected_coefficient_middle_family}); its exterior coefficients
    give the operators $B_f$ in (\ref{eq:selected_coefficient_decomposition}).
\end{proof}

\section{Trace contractions and interior slices}

The identities below give supporting contraction algebra for
\cite[Proposition~5.7]{Arad2016RigorousRG}. All matrix entries and scalar
coefficients are complex. Products retain the displayed order.

\begin{theorem}[Partition of an ordered product]
    \label{thm:trace_core_product_partition}
    \lean{TNLean.MPS.AreaLaw.orderedProduct_partition_interval}
    \leanok
    Let $l$ be a finite list in a monoid, and let $a,r$ be nonnegative
    integers. Write $l_{<a}$ for its prefix of length at most $a$,
    $l_{[a,a+r)}$ for the next at most $r$ entries, and $l_{\ge a+r}$ for
    the remaining suffix. Then
    \begin{align}
        \prod l=(\prod l_{<a})(\prod l_{[a,a+r)})(\prod l_{\ge a+r}).
        \label{eq:trace_core_product_partition}
    \end{align}
    These lists are truncated at the length of $l$, so no bound on $a+r$
    is required.
\end{theorem}
\begin{proof}\leanok
    Split the list after its first $a$ entries, split the remaining list
    after its first $r$ entries, and use associativity.
\end{proof}

\begin{theorem}[An open representation of a trace-closed chain]
    \label{thm:trace_core_open_square}
    \lean{TNLean.MPS.AreaLaw.isOpenMPS_trace_closed_chain}
    \leanok
    Let $L,q,D$ be nonnegative integers, let $C$ be a $D\times D$ matrix,
    and let $A_v(i)$ be a $D\times D$ matrix for each site $0\le v<L$
    and physical letter $0\le i<q$. The coefficient function
    \begin{align}
        \psi(w)=\operatorname{tr}\left(C\prod_{v=0}^{L-1}A_v(w_v)\right)
        \label{eq:trace_core_closed_coefficient}
    \end{align}
    has an open-boundary representation with internal bond dimensions at
    most $D^2$. At length zero its coefficient $\operatorname{tr}C$ is
    represented by the scalar multiplying the empty tensor product.
    Zero physical or virtual dimension is allowed.
\end{theorem}
\begin{proof}\leanok
    \uses{thm:chain_mps_of_cut_rank}
    At an internal cut, write the ordered matrix product as $P(x)Q(y)$.
    The coefficient matrix factors through pairs of virtual indices:
    \begin{align}
        \operatorname{tr}(CP(x)Q(y))
        =\sum_{i,j=0}^{D-1}(CP(x))_{ij}Q(y)_{ji}.
        \label{eq:trace_core_cut_factorization}
    \end{align}
    Every internal cut therefore has rank at most $D^2$. The
    characterization in Lemma~\ref{thm:chain_mps_of_cut_rank} gives
    the assertion.
\end{proof}

\begin{definition}[Interior matrix slices]
    \label{def:trace_core_matrix_slice}
    \lean{MPOTensor.interiorMatrixSlice}
    \leanok
    Let $A_v(i,j)$ be the supplied $B\times B$ site matrices of an operator
    chain of length $L$ and physical dimension $q$. For virtual indices
    $a,b$, its exposed interior slice is the operator $N_{ab}$ with entries
    \begin{align}
        N_{ab}(\sigma,\tau)
        =\left(\prod_{v=0}^{L-1}A_v(\sigma_v,\tau_v)\right)_{ab}.
        \label{eq:trace_core_matrix_slice}
    \end{align}
\end{definition}

\begin{definition}[A joint closing matrix for scalar word coefficients]
    \label{def:trace_core_word_closing}
    \lean{MPOTensor.stackClosingMatrix}
    \leanok
    Let $B,p$ be nonnegative integers. Encode a word
    $u\in\{0,\ldots,B-1\}^p$ by
    $\iota_p(u)=\sum_{r=0}^{p-1}B^{p-1-r}u_r$, so its last letter is the
    least significant digit. The empty word has index zero in a
    one-dimensional space, including when $B=0$.
    For a scalar function $c$ on words of pairs of virtual indices, define
    the $B^p\times B^p$ matrix $C(c)$ by
    \begin{align}
        C(c)_{\iota_p(v),\iota_p(u)}
        =c\bigl((u_r,v_r)_{r=0}^{p-1}\bigr).
        \label{eq:trace_core_word_closing}
    \end{align}
    This matrix depends only on $c$, $B$ and $p$.
\end{definition}

\begin{theorem}[Ordered slice words as a common stacked contraction]
    \label{thm:trace_core_slice_words}
    \lean{MPOTensor.sum_word_slices_eq_trace_stackedTensor}
    \leanok
    \uses{def:trace_core_matrix_slice,def:trace_core_word_closing}
    For the supplied site matrices $A_v$, let $A_v^{[p]}$ be their
    $p$-fold sitewise stack: the zero-fold stack is the identity tensor,
    and each successive layer contracts the intermediate physical index
    in the displayed operator order. Then, for all physical
    configurations $\sigma,\tau$,
    \begin{align}
         & \sum_{w=((a_r,b_r))_{r=0}^{p-1}}
        c(w)(N_{a_0b_0}\cdots N_{a_{p-1}b_{p-1}})(\sigma,\tau)\notag \\
         & \qquad=\operatorname{tr}\left(C(c)
        \prod_{v=0}^{L-1}A_v^{[p]}(\sigma_v,\tau_v)\right).
        \label{eq:trace_core_slice_words}
    \end{align}
    The same stacked site tensors serve every scalar coefficient function
    $c$. The closing matrix is chosen independently of the interior
    tensors. All lengths and dimensions may vanish whenever the displayed
    configurations exist.
\end{theorem}
\begin{proof}\leanok
    Fix endpoint words $x,y\in\{0,\ldots,B-1\}^{p}$, and write
    \begin{align*}
        P_p(x,y)(\sigma,\tau) & =(N_{x_0y_0}\cdots N_{x_{p-1}y_{p-1}})(\sigma,\tau), \\
        Q_p(x,y)(\sigma,\tau) & =\left(\prod_{v=0}^{L-1}A_v^{[p]}(\sigma_v,\tau_v)
        \right)_{\iota_p(x),\iota_p(y)}.
    \end{align*}
    We show $P_p(x,y)=Q_p(x,y)$ by induction on $p$. For $p=0$ both sides
    are the identity operator. For words $x,y$ of length $p+1$, let
    $\bar x,\bar y$ be their first $p$ letters, so that
    $\iota_{p+1}(x)=B\iota_p(\bar x)+x_p$. The operator product and the
    layer $A_v^{[p+1]}(\sigma_v,\tau_v)=\sum_{\rho_v}A_v^{[p]}(\sigma_v,\rho_v)
        \otimes A_v(\rho_v,\tau_v)$ give
    \begin{align}
        P_{p+1}(x,y)(\sigma,\tau)
         & =\sum_{\rho}P_p(\bar x,\bar y)(\sigma,\rho)\,
        N_{x_py_p}(\rho,\tau),\notag                     \\
        Q_{p+1}(x,y)(\sigma,\tau)
         & =\sum_{\rho}Q_p(\bar x,\bar y)(\sigma,\rho)
        \left(\prod_{v=0}^{L-1}A_v(\rho_v,\tau_v)\right)_{x_py_p},
        \label{eq:trace_core_layer_recursion}
    \end{align}
    where $\rho$ ranges over the physical configurations of the chain.
    By (\ref{eq:trace_core_matrix_slice}) the last factors in
    (\ref{eq:trace_core_layer_recursion}) agree, so the
    induction hypothesis gives $P_{p+1}(x,y)=Q_{p+1}(x,y)$. Separate each word of
    pairs into its two endpoint words and change to the encoded indices.
    The reversal of the row and column indices in
    (\ref{eq:trace_core_word_closing}) gives exactly the trace pairing in
    (\ref{eq:trace_core_slice_words}).
\end{proof}

\section{Retained matrices and finite families of ordered cores}

The identities below give supporting algebra for
\cite[Proposition~5.7]{Arad2016RigorousRG}.
All products of site matrices retain their natural order.

\begin{definition}[An interval and its single retained complement]
    \label{def:retained_interval_coordinates}
    \lean{TNLean.MPS.AreaLaw.retainedIntervalEquiv,MPOTensor.actualAffineRetainedConfig}
    Let $T,L,a$ be nonnegative integers with $a+L\le T$, and set
    $J=[a,a+L)$ in the chain of $T$ sites. For any set $X$, restriction gives
    an equivalence
    \begin{align}
        e:X^{J^c}\times X^{\{0,\ldots,L-1\}} & \longrightarrow
        X^{\{0,\ldots,T-1\}},\notag                            \\
        e(\eta,x)_v                          & =
        \begin{cases}x_{v-a},&a\le v<a+L,\\ \eta_v,&v\notin J.\end{cases}
        \label{eq:retained_interval_coordinates}
    \end{align}
    Its inverse restricts a configuration to $J^c$ and to the sites $a+r$,
    $0\le r<L$. For local alphabet $\{0,\ldots,q-1\}$, denote the retained
    configuration set by $\mathcal R=\{0,\ldots,q-1\}^{J^c}$.
    The complement is a single set: each site outside $J$ occurs once.
\end{definition}

\begin{definition}[The actual retained closing matrix]
    \label{def:actual_retained_interval_closing}
    \lean{MPOTensor.actualRetainedIntervalClosing}
    \uses{def:retained_interval_coordinates}
    Let $A_v(i,j)\in\C^{B\times B}$ be supplied site-dependent operator
    tensors on the full chain, and let $C\in\C^{B\times B}$ be supplied.
    For $\eta,\theta\in\mathcal R$, define
    \begin{align}
        P(\eta,\theta) & =\prod_{v=0}^{a-1}A_v(\eta_v,\theta_v),         &
        S(\eta,\theta) & =\prod_{v=a+L}^{T-1}A_v(\eta_v,\theta_v),\notag   \\
        F(\eta,\theta) & =S(\eta,\theta)\,C\,P(\eta,\theta).
        \label{eq:actual_retained_interval_closing}
    \end{align}
    Empty products are identity matrices. The matrix $F(\eta,\theta)$ is
    chosen before any physical coordinates inside $J$.
\end{definition}


\begin{theorem}[Restriction of a trace contraction to an interval]
    \label{thm:actual_retained_interval_trace}
    \lean{MPOTensor.trace_actual_retainedInterval_eq}
    \uses{def:retained_interval_coordinates,def:actual_retained_interval_closing,
        def:trace_chain_operator}
    Use $a+L\le T$, the full-chain tensors $A_v$, the closing matrix $C$,
    the coordinate equivalence $e$, and the fixed retained configurations
    $\eta,\theta$ of (\ref{eq:actual_retained_interval_closing}).
    Put $A^{\mathrm{int}}_r=A_{a+r}$ for $0\le r<L$.
    For all interior configurations $\sigma,\tau$,
    \begin{align}
        \mathcal T(A,C)(e(\eta,\sigma),e(\theta,\tau))
        &=\mathcal T(A^{\mathrm{int}},F(\eta,\theta))(\sigma,\tau).
        \label{eq:actual_retained_interval_trace}
    \end{align}
    Empty intervals and zero dimensions are allowed whenever the supplied
    configurations exist.
\end{theorem}
\begin{proof}
    Split the ordered product into prefix, interval and suffix.
    Cyclicity gives $\operatorname{tr}(CPXS)=\operatorname{tr}(SCPX)$,
    and $F=SCP$ by (\ref{eq:actual_retained_interval_closing}).
\end{proof}

\begin{theorem}[Insertion of fixed configurations at the same trace dimension]
    \label{thm:retained_single_trace_insert}
    \lean{MPOTensor.exists_traceChainOperator_retainedSingle_insert}
    \uses{def:retained_interval_coordinates,def:trace_chain_operator}
    Let $T,L,q,\chi,a$ be nonnegative integers with $a+L\le T$.
    Put $J=[a,a+L)$ and supply site tensors $A_r$ on the $L$ interior sites,
    of virtual dimension $\chi$, and a closing matrix $C$.
    Fix row and column configurations $\eta,\theta$ on $J^c$.
    There are tensors $A'_v$ on the full $T$ sites, of the same dimension,
    such that for all full-chain configurations $\sigma,\tau$,
    \begin{align}
        \mathcal T(A',C)(\sigma,\tau)
        &=\mathbf1_{\sigma|_{J^c}=\eta}\mathbf1_{\tau|_{J^c}=\theta}
        \mathcal T(A,C)(\sigma|_J,\tau|_J).
        \label{eq:retained_single_trace_insert}
    \end{align}
    The restrictions to $J$ use its increasing enumeration. Both $L$ and
    $T$ may vanish; no positivity assumption on $q$ or $\chi$ is imposed.
\end{theorem}
\begin{proof}
    Use $A'_v=A_{v-a}$ inside $J$ and
    $A'_v(i,j)=\mathbf1_{i=\eta_v}\mathbf1_{j=\theta_v}\Id_\chi$ outside $J$.
    If the retained configurations match, the extra letters are identity
    matrices and the trace is unchanged. Otherwise one letter is zero.
\end{proof}

\begin{theorem}[A trace contraction after restriction and fixed extension]
    \label{thm:trace_interval_slice_configuration_extend}
    \lean{MPOTensor.exists_traceChainOperator_intervalSlice_extend_by_configuration}
    \uses{thm:actual_retained_interval_trace,thm:retained_single_trace_insert}
    Let $T,L,N,q,\chi,a,b$ be nonnegative integers with $a+L\le T$ and
    $b+L\le N$. Supply old tensors $A_v$ of dimension $\chi$, a closing
    matrix $C$, and full old configurations $f,g$. For an interior
    configuration $z$, write $\iota_a(f,z)$ for the old configuration
    equal to $z_{v-a}$ on $[a,a+L)$ and to $f_v$ elsewhere.
    Fix new outside configurations $\eta,\theta$ on
    $J'^c$, where $J'=[b,b+L)$ in the new chain of $N$ sites.
    There are a matrix $C'$ and tensors $A'_v$ of the same virtual
    dimension $\chi$ such that
    \begin{align}
        \mathcal T(A',C')(\sigma,\tau)
        &=\mathbf1_{\sigma|_{J'^c}=\eta}\mathbf1_{\tau|_{J'^c}=\theta}
        \mathcal T(A,C)\bigl(\iota_a(f,\sigma|_{J'}),
            \iota_a(g,\tau|_{J'})\bigr).
        \label{eq:trace_interval_slice_configuration_extend}
    \end{align}
    Restrictions to $J'$ use increasing coordinates. All chain lengths and
    dimensions may vanish, subject only to the stated supplied data.
\end{theorem}
\begin{proof}
    \uses{thm:actual_retained_interval_trace,thm:retained_single_trace_insert}
    Fix the old outside configurations to those of $f,g$ and apply
    Theorem~\ref{thm:actual_retained_interval_trace}, obtaining the restricted
    tensors and their suffix-closing-prefix matrix $C'$.
    Insert the new retained configurations by
    Theorem~\ref{thm:retained_single_trace_insert}, retaining $C'$ and $\chi$.
\end{proof}

\begin{definition}[The scalar-closed core at a retained matrix entry]
    \label{def:retained_stack_core}
    \lean{MPOTensor.retainedStackCore,MPOTensor.exists_retainedStackCore_traceChainOperator}
    \uses{def:core_word_coefficients,def:trace_core_word_closing}
    Let $m$ be a finite retained index set, let
    $R_{ab}\in\C^{m\times m}$ for $0\le a,b<B$, and let
    $s=((K_0,c_0),\ldots,(K_{h-1},c_{h-1}))$ be a list of affine steps.
    For $t\ge0$ and $i,j\in m$, let
    \begin{align}
        \gamma_t^{i,j}(w)=C_R(s,w)_{ij},
        \qquad w\in(\{0,\ldots,B-1\}^2)^t.
        \label{eq:retained_stack_scalar_coefficient}
    \end{align}
    The joint closing matrix $D_t^{i,j}$ is the closing matrix associated to
    these scalar word coefficients. In the stacking coordinates $\iota_t$
    on words of $t$ virtual labels, its entries are
    \begin{align}
        (D_t^{i,j})_{\iota_t(v),\iota_t(u)}
        =\gamma_t^{i,j}\bigl((u_r,v_r)_{r=0}^{t-1}\bigr).
        \label{eq:retained_stack_closing_entries}
    \end{align}
    For supplied interior tensors $A_v$, $0\le v<L$, let $A_v^{[t]}$ be
    their existing $t$-layer stack, of virtual dimension $B^t$. Define
    \begin{align}
        W_t^{i,j}(\sigma,\tau)=\operatorname{tr}\left(
        D_t^{i,j}\prod_{v=0}^{L-1}A_v^{[t]}(\sigma_v,\tau_v)\right).
        \label{eq:retained_stack_core}
    \end{align}
    The closing matrix depends on $R,s,t,i,j$, independently of the interior
    chain. In particular, no rank-one condition on that matrix is imposed.
\end{definition}

\begin{theorem}[A finite family in natural physical order]
    \label{thm:finite_natural_retained_core_family}
    \lean{MPOTensor.exists_finite_natural_retained_core_family}
    \uses{def:retained_interval_coordinates,def:retained_stack_core,
        def:trace_core_matrix_slice,thm:core_word_product,
        thm:trace_core_slice_words,thm:trace_core_open_square,
        thm:chain_fixed_configuration_extension}
    Fix $a+L\le T$, interior tensors $A_v$ of virtual dimension $B$, and
    a retained index set $m=\mathcal R$. Choose $R_{ab}\in\C^{m\times m}$
    and a list $s$ of $h$ affine steps on $m$.
    Let $N_{ab}(A)$ be the interior matrix slice obtained by fixing the two
    virtual endpoints to $a,b$. Put $n=(h+1)|m|^2$.
    There are degrees $t_r\in\{0,\ldots,h\}$, retained indices $i_r,j_r\in m$,
    and full-chain operators $V_r$, for $0\le r<n$, such that
    \begin{align}
        V_r & =\left(|i_r\rangle\langle j_r|\otimes
        W_{t_r}^{i_r,j_r}\right)^e,\notag           \\
        \left(\prod_{v=0}^{h-1}\left(K_v\otimes I+
        \sum_{a,b}(c_vR_{ab})\otimes N_{ab}(A)\right)\right)^e
            & =\sum_{r=0}^{n-1}V_r.
        \label{eq:finite_natural_retained_core_family}
    \end{align}
    The superscript denotes transport from retained/interior coordinates by
    the equivalence $e$. Each $V_r$, viewed as a vector on the paired physical
    alphabet, has an open-boundary MPS representation with bond bound
    $(B^{t_r})^2$. Empty chains, empty intervals, and zero physical or virtual
    dimensions are included whenever the stated retained data exist.
\end{theorem}
\begin{proof}
    \uses{def:retained_interval_coordinates,def:retained_stack_core,
        thm:core_word_product,thm:trace_core_slice_words,thm:trace_core_open_square,
        thm:chain_fixed_configuration_extension}
    Expand the ordered affine product over words of each length $0\le t\le h$.
    Resolve every retained coefficient on the matrix units of $m$. For fixed
    $t,i,j$, the scalar word coefficients are exactly
    (\ref{eq:retained_stack_scalar_coefficient}). Their joint closing matrix
    therefore gives the core (\ref{eq:retained_stack_core}). Enumerate the
    finite set $\{0,\ldots,h\}\times m\times m$ to obtain the literal family
    in (\ref{eq:finite_natural_retained_core_family}).

    A trace-closed chain of virtual dimension $B^t$ has open-boundary bond
    bound $(B^t)^2$. Inserting fixed retained configurations outside the
    interval gives bond bound $\max\{1,(B^t)^2\}$. If $B^t>0$, this is the
    stated square bound. If $B^t=0$, the trace defining the core vanishes;
    the full inserted operator is zero and has the square bound as well.
\end{proof}

\begin{theorem}[The affine product of the actual supplied trace operator]
    \label{thm:actual_natural_affine_product}
    \lean{MPOTensor.prod_actual_natural_affine_eq}
    \uses{def:retained_interval_coordinates,def:actual_retained_interval_closing}
    Fix $a+L\le T$. For supplied full-chain tensors $A^{\mathrm{full}}_v$
    of virtual dimension $B$ and cap $C\in\C^{B\times B}$, define
    \begin{align}
        H(\sigma,\tau)      & =\operatorname{tr}\left(
        C\prod_{v=0}^{T-1}A^{\mathrm{full}}_v(\sigma_v,\tau_v)\right),\notag \\
        A^{\mathrm{int}}_r  & =A^{\mathrm{full}}_{a+r}, &
        R_{ab}(\eta,\theta) & =F(\eta,\theta)_{ba}.
        \label{eq:actual_natural_affine_data}
    \end{align}
    For every ordered list $s=((K_v,c_v))_{v=0}^{h-1}$ on $\mathcal R$,
    \begin{align}
        \prod_{v=0}^{h-1}\left((K_v\otimes I)^e+c_vH\right)
        =\left(\prod_{v=0}^{h-1}\left(K_v\otimes I+
        \sum_{a,b}(c_vR_{ab})\otimes N_{ab}(A^{\mathrm{int}})\right)\right)^e.
        \label{eq:actual_natural_affine_product}
    \end{align}
    Here $F$ is exactly the suffix-cap-prefix matrix of
    (\ref{eq:actual_retained_interval_closing}). The identity concerns the
    supplied trace operator $H$; no identification with another Hamiltonian
    is assumed.
\end{theorem}
\begin{proof}
    \uses{def:retained_interval_coordinates,def:actual_retained_interval_closing}
    Split the ordered full-chain product into prefix, interval and suffix.
    Cyclicity of the trace moves the suffix before the cap, giving $F=SCP$.
    Expanding the trace over its two virtual indices gives
    $H^{{e}^{-1}}=\sum_{a,b}R_{ab}\otimes N_{ab}(A^{\mathrm{int}})$;
    the reversed indices in $F_{ba}$ come from the trace pairing.
    Transport preserves sums, scalar multiplication and ordered products.
    Applying it to each affine factor proves
    (\ref{eq:actual_natural_affine_product}).
\end{proof}

\section{Natural coordinates for physical interval operators}

The following identities give the coordinate and placement steps supporting
\cite[Proposition~5.7]{Arad2016RigorousRG}. All endpoints and dimensions below
are nonnegative integers. Intervals are half-open and intersected with the
finite chain. Subtraction of natural numbers is truncated at zero.

\begin{definition}[Increasing coordinates on an interval]
    \label{def:physical_interval_increasing_coordinates}
    \lean{TNLean.MPS.AreaLaw.intervalShiftOrderIso}
    For $a\le b\le T$, the increasing coordinate identification is the
    order isomorphism
    \begin{align}
        e_{a,b}:\{0,\ldots,b-a-1\} & \longrightarrow[a,b), &
        e_{a,b}(i)                 & =a+i.
        \label{eq:physical_interval_increasing_coordinates}
    \end{align}
    Its inverse sends a site $x$ to $x-a$. Empty intervals are included.
\end{definition}

\begin{theorem}[A subinterval in increasing coordinates]
    \label{thm:physical_interval_subinterval_coordinates}
    \lean{TNLean.MPS.AreaLaw.intervalShiftOrderIso_mem_subinterval}
    \uses{def:physical_interval_increasing_coordinates}
    Suppose $c\le d\le T$ and $c\le u$. For each $i<d-c$,
    \begin{align}
        e_{c,d}(i)\in[u,v)
        \quad\Longleftrightarrow\quad
        u-c\le i<(u-c)+(v-u).
        \label{eq:physical_interval_subinterval_coordinates}
    \end{align}
    No ordering of $u,v$ is required: if $v<u$, both conditions are false.
\end{theorem}
\begin{proof}
    Substitute $e_{c,d}(i)=c+i$ and compare the endpoint inequalities.
\end{proof}

\begin{theorem}[The length of the natural subinterval]
    \label{thm:physical_interval_subinterval_length}
    \lean{TNLean.MPS.AreaLaw.intervalShift_subinterval_length_le}
    If $c\le u\le d$ and $v\le d$, then
    \begin{align}
        (u-c)+(v-u)\le d-c.
        \label{eq:physical_interval_subinterval_length}
    \end{align}
    This includes $v<u$, when the subinterval length is zero.
\end{theorem}
\begin{proof}
    If $u\le v$, the left side equals $v-c$. If $v<u$, it equals $u-c$.
    Each is at most $d-c$.
\end{proof}

\begin{definition}[Coordinates on the single retained complement]
    \label{def:physical_interval_complement_coordinates}
    \lean{TNLean.MPS.AreaLaw.intervalComplementEquiv}
    \uses{def:physical_interval_increasing_coordinates,
        thm:physical_interval_subinterval_coordinates}
    For $c\le d\le T$ and $c\le u$, restriction of $e_{c,d}$ gives a
    bijection
    \begin{align}
        \{i<d-c:i\notin[u-c,(u-c)+(v-u))\}
        \longrightarrow
        [c,d)\setminus
        [u,v).
        \label{eq:physical_interval_complement_coordinates}
    \end{align}
    The retained sites form one complement; a site belonging to overlapping
    endpoint regions occurs only once.
\end{definition}

\begin{definition}[Identity extension of the supplied core tensors]
    \label{def:physical_interval_identity_extension}
    \lean{MPOTensor.identityExtendInterval}
    Let $A_r(i,j)$, $0\le r<L$, be the supplied $B\times B$ site matrices,
    and suppose $a+L\le T$. Their extension to the $T$-site chain is
    \begin{align}
        \widehat A_v(i,j)=
        \begin{cases}
            A_{v-a}(i,j),   & a\le v<a+L,     \\
            \delta_{ij}I_B, & v\notin[a,a+L).
        \end{cases}
        \label{eq:physical_interval_identity_extension}
    \end{align}
    The virtual space remains the supplied $B$-dimensional space.
\end{definition}

\begin{theorem}[Reading a nested physical interval in natural coordinates]
    \label{thm:physical_interval_nested_trace}
    \lean{MPOTensor.trace_identityExtendInterval_nested_natural}
    \uses{def:physical_interval_increasing_coordinates,
        def:physical_interval_identity_extension}
    Suppose $c\le a\le b\le d\le T$. Let $A_r(i,j)$ be the supplied core
    tensors on $b-a$ sites, with scalar closing matrix $C$ on the same
    $B$-dimensional virtual space. Let $H_K$ be their trace-form operator
    placed on $K=[a,b)$ using $e_{a,b}$, and extend it canonically to
    $E=[c,d)$. Reading this operator in the increasing coordinates of $E$
    gives the matrix with entries
    \begin{align}
        \operatorname{tr}\left(C\prod_{v=0}^{d-c-1}
        \widehat A_v(\sigma_v,\tau_v)\right),
        \label{eq:physical_interval_nested_trace}
    \end{align}
    where the identity extension starts at $a-c$ inside the natural chain
    of length $d-c$. The closing matrix and virtual space are unchanged.
    Empty intervals and zero physical or virtual dimension are included.
\end{theorem}
\begin{proof}
    When the two physical configurations agree outside the core interval,
    all exterior letters are identities, and the ordered product is the
    original core product. Otherwise one exterior letter is zero. This
    is exactly the canonical extension of the core operator. Composition
    of the increasing interval inclusions gives the stated physical
    placement, and reading along the inverse increasing coordinates
    recovers its natural matrix.
\end{proof}

\begin{theorem}[A physical complement operator as a retained matrix]
    \label{thm:physical_interval_retained_matrix}
    \lean{TNLean.MPS.AreaLaw.exists_retainedMatrix_of_physical_localLift}
    Let $V$ be a finite physical site set and $e:\{0,\ldots,T-1\}\to V$
    a bijection. Fix $a,L$ with $a+L\le T$, and suppose the physical set
    $C_{\rm phys}$ consists exactly of the images under $e$ of the sites
    outside $[a,a+L)$. Let $K_{\rm phys}$ be any operator on its
    configuration space, with local dimension $q$.
    There is a matrix $K_{\rm ret}$ on configurations of the single
    natural complement such that the canonical extension of $K_{\rm phys}$,
    first read along $e$ and then split into retained and interior
    configurations, is
    \begin{align}
        K_{\rm ret}\otimes I_{q^L}.
        \label{eq:physical_interval_retained_matrix}
    \end{align}
    The retained matrix is obtained from the supplied regional matrix
    by this coordinate identification. No interior operator is required.
\end{theorem}
\begin{proof}
    \uses{def:retained_interval_coordinates}
    Relabel the canonical extension along $e$ using covariance of the
    regional lift. The resulting natural complement matrix is reindexed
    by the complement configuration bijection. Splitting a full
    configuration into its complement and interval parts leaves the
    complement entry unchanged and contributes an identity entry on the
    interval.
\end{proof}

\begin{theorem}[A trace-chain representation on a supported interval]
    \label{thm:supported_interval_trace_chain}
    \lean{TNLean.MPS.AreaLaw.exists_traceChainOperator_of_supported_interval}
    \uses{def:chain_model,def:trace_chain_operator,
        def:physical_interval_increasing_coordinates}
    Let $0\le c\le d\le L$ and let $A$ be supported on the interval
    $J=\{c,\ldots,d-1\}$. There are a finite bond dimension, site tensors
    on a chain of $d-c$ sites and a closing matrix whose trace contraction,
    transported by $i\mapsto c+i$ and extended by the identity outside $J$,
    equals $A$. Empty intervals and zero local dimension are included.
\end{theorem}
\begin{proof}
    \uses{thm:finite_matrix_trace_chain,def:physical_interval_increasing_coordinates}
    The support hypothesis supplies a matrix on the interval coordinates.
    Pull it back along the increasing identification, represent that finite
    matrix by a trace chain, and transport the representation to its sites.
    This supplies no efficient bond bound.
\end{proof}

\begin{theorem}[Inward effective cuts with a bounded collar]
    \label{thm:inward_effective_cuts}
    \lean{TNLean.MPS.AreaLaw.inward_effective_cuts}
    \uses{def:chain_model}
    Let $\ell\ge1$, $0\le p\le r\le L$ and $r-p>4\ell$.
    Put $p'=\max(p,\ell)$, $r'=\min(r,L-\ell)$,
    $a=p'-\ell$ and $b=r'-\ell-1$. Then
    \begin{align}
        a+\ell  & =p',               & b+\ell+1 & =r',\notag        \\
        a+2\ell & \le b,             & b+2\ell  & \le L-1,\notag    \\
        p       & \le p'\le r'\le r, & r'-p'    & \ge2\ell+1,\notag \\
        p'-p    & \le\ell,           & r-r'     & \le\ell.
    \end{align}
    Also $(p'-p)+(r-r')\le2\ell$ and, for
    $J_0=\{p',\ldots,r'-1\}$ and $J=\{p,\ldots,r-1\}$,
    $J_0\subseteq J$ and $|J\setminus J_0|\le2\ell$.
\end{theorem}
\begin{proof}
    The long-interval assumption allows moving each endpoint inward by at
    most $\ell$ while retaining length at least $2\ell+1$.
    Substitute the definitions in the endpoint inequalities. The removed
    collar has $(p'-p)+(r-r')$ sites, giving its cardinal bound.
\end{proof}

\section{Middle families on actual regions}

The following regional identities give supporting algebra for changing
the two cuts in
\cite[Proposition~5.7]{Arad2016RigorousRG}.
The matrix-factor exchange is the existing construction in
Definition~\ref{def:chain_middle_factor_exchange} and
Theorem~\ref{thm:chain_middle_factor_exchange}; its separate cardinality
identity is Theorem~\ref{thm:chain_middle_factor_exchange_card}.
No ordering of physical sites or matrix product representation is asserted
by these regional identities.

\begin{theorem}[From numbered middle coordinates to the actual middle region]
    \label{thm:regional_middle_from_numbered_three_regions}
    \lean{TNLean.MPS.AreaLaw.exists_regional_exterior_of_natural_threeSplit_decomposition}
    \uses{def:selected_exterior_middle_equiv}
    Let $\Omega$ be a finite site set, let $T,E\subseteq\Omega$ be disjoint,
    and put $R=\Omega\setminus(T\cup E)$. Fix a local alphabet
    $\{0,\ldots,q-1\}$, a semiring $\mathbb S$, a finite index set $I$,
    and a bijection $e_E:\{0,\ldots,N-1\}\to E$.
    Suppose a full matrix $X$ has a decomposition in the coordinates
    $((\sigma_T,\sigma_R),x)$, where $x$ is a configuration on
    $\{0,\ldots,N-1\}$, of the form
    \begin{align}
        X\bigl(g^{-1}((l,x\circ e_E^{-1}),r),
        g^{-1}((l',y\circ e_E^{-1}),r')\bigr)
        =\sum_{t\in I}B_t((l,r),(l',r'))V_t(x,y).
        \label{eq:regional_middle_numbered_decomposition}
    \end{align}
    Here $g$ restricts a full configuration to $T,E,R$, in that order.
    Then there are matrices $M_t$ on configurations of $E^c$ such that
    \begin{align}
        X(\sigma,\tau)=\sum_{t\in I}
        M_t(\sigma|_{E^c},\tau|_{E^c})
        V_t\bigl((\sigma|_E)\circ e_E,(\tau|_E)\circ e_E\bigr).
        \label{eq:regional_middle_actual_decomposition}
    \end{align}
    The middle matrices are the prescribed $V_t$ transported by $e_E$; only
    the exterior coefficients change. The bijection need not carry any order.
\end{theorem}
\begin{proof}
    The two inverse identities for $e_E$ first turn the given decomposition
    into one with actual $E$ configurations. Restriction identifies the actual
    exterior configuration on $E^c$ with its pair of restrictions to $T,R$.
    Pull each $B_t$ back along this latter identification. The two coordinate
    descriptions of a full configuration then agree entry by entry, giving
    (\ref{eq:regional_middle_actual_decomposition}).
\end{proof}

\begin{definition}[The old and new middle configuration splittings]
    \label{def:regional_middle_configuration_splittings}
    \lean{FiniteProduct.changingOldMiddleConfigEquiv,FiniteProduct.changingNewMiddleConfigEquiv}
    Let $\Omega$ be finite, let $\beta_v$ be any set for each site, and write
    $\mathcal X_S=\prod_{v\in S}\beta_v$. For regions $E,J\subseteq\Omega$, put
    \begin{align}
        A=E\setminus J,\qquad B=J\setminus E,\qquad
        C=E\cap J,\qquad O=(E\cup J)^c.
        \label{eq:regional_middle_four_regions}
    \end{align}
    The old and new middle splittings are the restriction equivalences
    \begin{align}
        \phi_E:\mathcal X_E & \longrightarrow\mathcal X_A\times\mathcal X_C, &
        \phi_J:\mathcal X_J & \longrightarrow\mathcal X_B\times\mathcal X_C.
        \label{eq:regional_middle_configuration_splittings}
    \end{align}
    Both use the same literal intersection $C$. No nonemptiness or finiteness
    of the local sets $\beta_v$ is required for these equivalences.
\end{definition}

\begin{definition}[The prescribed new regional middle family]
    \label{def:regional_middle_transferred_family}
    \lean{FiniteProduct.changingRegionalMiddleFamily}
    \uses{def:regional_middle_configuration_splittings,def:chain_middle_factor_exchange}
    Let $\mathbb S$ be a semiring, let $I$ be any set, and let $V_t$ be
    $\mathbb S$-valued matrices on $\mathcal X_E$, indexed by $t\in I$.
    For $r=(t,(a_0,a_1),(b_0,b_1))$ in
    \begin{align}
        I_{\mathrm{new}}=I\times\mathcal X_A^2\times\mathcal X_B^2,
        \label{eq:regional_middle_new_index}
    \end{align}
    write $V_t$ in the coordinates $\phi_E$ and denote its partial matrix
    element at $a_0,a_1$ by $(V_t^{\phi_E})_{a_0,a_1}$.
    Define the matrix $\widetilde V_r$ on $\mathcal X_J$ by
    \begin{align}
        \widetilde V_r(y,y')=
        \left(E_{b_0b_1}\otimes (V_t^{\phi_E})_{a_0,a_1}\right)
        (\phi_J y,\phi_J y').
        \label{eq:regional_middle_transferred_family}
    \end{align}
    Here $E_{b_0b_1}$ is the standard matrix unit with entry $1$ at $(b_0,b_1)$.
    The family depends only on the old middle matrices and the two regions,
    before any exterior coefficient matrices are supplied.
\end{definition}

\begin{theorem}[Transfer of a fixed middle family between regions]
    \label{thm:regional_middle_exterior_transfer}
    \lean{FiniteProduct.exists_changingRegional_exterior_family}
    \uses{def:regional_middle_configuration_splittings,
        def:regional_middle_transferred_family,thm:chain_middle_factor_exchange}
    Assume now that $I$ and every local set $\beta_v$ are finite.
    For every family $M_t$ of $\mathbb S$-valued matrices on
    $\mathcal X_{E^c}$, there are matrices $M^{\mathrm{new}}_r$ on
    $\mathcal X_{J^c}$ such that, for all full configurations $\sigma,\tau$,
    \begin{align}
        \sum_{t\in I}M_t(\sigma|_{E^c},\tau|_{E^c})
        V_t(\sigma|_E,\tau|_E)
        =\sum_{r\in I_{\mathrm{new}}}
        M^{\mathrm{new}}_r(\sigma|_{J^c},\tau|_{J^c})
        \widetilde V_r(\sigma|_J,\tau|_J).
        \label{eq:regional_middle_exterior_transfer}
    \end{align}
    The family $\widetilde V$ is the one already prescribed above and is
    independent of $M$. Empty regions and empty local sets are included.
\end{theorem}
\begin{proof}
    \uses{def:regional_middle_configuration_splittings,
        def:regional_middle_transferred_family,thm:chain_middle_factor_exchange}
    The old split into the four regions gives $((\mathcal X_B,\mathcal X_O),(\mathcal X_A,\mathcal X_C))$.
    The new split gives
    $((\mathcal X_A,\mathcal X_O),(\mathcal X_B,\mathcal X_C))$.
    Both descriptions are literal restrictions of the same full configuration,
    so the existing matrix-factor exchange carries one to the other.

    Apply that exchange to the old sum of Kronecker products. Its new middle
    factors are precisely (\ref{eq:regional_middle_transferred_family}).
    Pull the new exterior factors back to configurations on $J^c$ and the new
    middle factors back to configurations on $J$. Evaluating at the two full
    configurations proves (\ref{eq:regional_middle_exterior_transfer}), with
    the multiplication order of the semiring coefficients unchanged.
\end{proof}

\section{A common natural family for selected coefficients}

The following results provide supporting algebra and coarse counts for
\cite[Proposition~5.7]{Arad2016RigorousRG}.
The number of family members and their open-boundary bond dimensions are
separate quantities. No sharp source bond or running-time conclusion is
included.

\begin{theorem}[Natural middle operators chosen before the exterior pair]
    \label{thm:selected_common_natural_mpo_family}
    \lean{TNLean.MPS.AreaLaw.exists_common_physical_twoMarkerCoefficient_mpo_family}
    \uses{thm:selected_coefficient_prescribed_middle,
        thm:selected_middle_fixed_physical_data,thm:finite_natural_retained_core_family,
        thm:actual_natural_affine_product,thm:physical_interval_nested_trace,
        thm:physical_interval_retained_matrix,def:two_marker_middle_index}
    Let the chain have $L$ sites with local dimension $q\ge0$.
    Let $a,b,\ell$ be nonnegative integers satisfying
    $a+2\ell\le b$ and $b+2\ell\le L-1$, and choose $0\le p,s<2\ell$.
    Use the regions
    \begin{align}
        T & ={[0,a+p+1)}, & E &
        ={[a+p+1,b+s+1)},\notag                 \\
        R & =(T\cup E)^c, & Z & =[a+2\ell+1,b),
        \label{eq:selected_common_natural_regions}
    \end{align}
    with all intervals clipped to the chain. Set $N=|E|$, and identify $E$
    with its increasing list of sites. Let $\rho$ be the number of
    configurations on $E\setminus Z$.

    Suppose the supplied window operators $W^L_r,W^R_r$ are supported on
    $[a+r,a+r+2)$ and
    $[b+r,b+r+2)$, respectively.
    Supply site-dependent core tensors of virtual dimension $B$ on
    $[a+2\ell,b+1)$ and a closing matrix. Let $P_M$ be their trace contraction,
    canonically placed on the full chain.
    Use labels $L_r,R_r,B_0,B_1,B_2$, and put
    $O=\{L_r,R_r,B_0,B_1,B_2\}\setminus\{L_p,R_s\}$.
    Fix complex ordinary weights $z_t$ and nonnegative integers $k,i,j$.
    Write $\mathcal W_{k,i,j}$ for the two-marker middle index set with
    $q^2$ labels for each marker family, as in
    Definition~\ref{def:two_marker_middle_index}.

    There exist a finite family of natural-interval operators $V_r$,
    $0\le r<n$, and degrees $0\le d_r\le k$, such that
    \begin{align}
        n & \le |\mathcal W_{k,i,j}|(k+1)\rho^2,
        \label{eq:selected_common_natural_count}
    \end{align}
    and every $V_r$, on the paired physical alphabet, has an open-boundary
    MPS representation of bond at most $(B^{d_r})^2$ in increasing site order.
    These choices precede the following exterior data.

    For every supported exterior pair $P_L,P_R$ on $[0,a+1)$ and
    $[b+2\ell,L)$, respectively, form the full family
    $F_{L_r}=W^L_r$, $F_{R_r}=W^R_r$,
    $F_{B_0}=P_L$, $F_{B_1}=P_M$, $F_{B_2}=P_R$, and set
    \begin{align}
        D_z & =\sum_{t\in O}z_tF_t.
        \label{eq:selected_common_natural_ordinary}
    \end{align}
    For every complex scalar family $c_0,\ldots,c_k$, there are joint exterior
    operators $Q_r$ on $\mathcal H_T\otimes\mathcal H_R$ such that
    \begin{align}
        \left(\sum_{h=0}^k c_h\,[u^iv^j]
        (D_z+uW^L_p+vW^R_s)^h\right)^{\mathrm{nat}}
         & =\sum_{r=0}^{n-1}Q_r\otimes V_r.
        \label{eq:selected_common_natural_decomposition}
    \end{align}
    The superscript first groups the two exterior regions and then reads the
    middle in its increasing coordinates. The formal variables $u,v$ commute
    with the matrix coefficients. The two selected crossings are unweighted;
    only the ordinary sum carries $z$.
    Zero physical or virtual dimensions, degree zero, and impossible crossing
    frequencies are included. No factorization of each joint exterior $Q_r$
    into separate left and right operators is asserted.
\end{theorem}
\begin{proof}
    \uses{thm:selected_coefficient_prescribed_middle,
        thm:selected_middle_fixed_physical_data,thm:finite_natural_retained_core_family,
        thm:actual_natural_affine_product,thm:physical_interval_nested_trace,
        thm:physical_interval_retained_matrix,def:two_marker_middle_index,
        thm:two_marker_prescribed_span}
    The prescribed-middle coefficient decomposition first gives one ordered
    middle factor for every index $w\in\mathcal W_{k,i,j}$, before the exterior
    pair or the scalars $c_h$ are chosen. Its ordinary factor has the form
    $z_{B_1}M_E+K_z$, with $K_z$ supported on $E\setminus Z$; its two marker
    families are also supported there.

    Represent those supported factors on the single retained complement.
    Replace each ordinary run by its repeated affine steps, with scalar
    $z_{B_1}$ on each ordinary step and scalar zero on each marker step.
    The resulting ordered list has length $h_w\le k$ and evaluates to the
    same middle word. The canonical interval and retained-coordinate
    identities identify its central trace operator with the supplied core.
    The finite natural family theorem expresses that word as a sum of exactly
    $(h_w+1)\rho^2$ natural operators, each with its corresponding square bond.

    Take the disjoint union over $w$ and enumerate it. Its size is at most
    $|\mathcal W_{k,i,j}|(k+1)\rho^2$, proving
    (\ref{eq:selected_common_natural_count}). Every prescribed middle word lies
    in the span of this one family. The prescribed-spanning-family
    factorization then supplies the exterior coefficients in
    (\ref{eq:selected_common_natural_decomposition}), without changing the
    chosen middle operators.
\end{proof}

\begin{theorem}[A polynomial-base bound for the selected family count]
    \label{thm:selected_common_natural_cardinality}
    \lean{TNLean.MPS.AreaLaw.actual_common_physical_mpo_card_le}
    \uses{thm:selected_middle_collar_cardinality,thm:two_marker_middle_cardinality,
        def:physical_interval_complement_coordinates}
    Retain the geometric data and positions above, and suppose $q\ge1$ and
    $k\ge1$. For arbitrary nonnegative $i,j,n$, the numerical implication
    \begin{align}
        n\le |\mathcal W_{k,i,j}|(k+1)\rho^2
        \quad\Longrightarrow\quad
        n\le\bigl(q(k+1)\bigr)^{8\ell+2(i+j)+2}
        \label{eq:selected_common_natural_cardinality}
    \end{align}
    holds. No low-frequency restriction is required.
\end{theorem}
\begin{proof}
    \uses{thm:selected_middle_collar_cardinality,thm:two_marker_middle_cardinality,
        def:physical_interval_complement_coordinates}
    The collar bound gives $\rho\le q^{4\ell}$, and hence
    $\rho^2\le q^{8\ell}$. With $m=i+j$ and $v=k+1$, the middle index count is
    at most $2^m q^{2m}v^{m+1}$. Since $v\ge2$, it is at most
    $q^{2m}v^{2m+1}$. Multiplication by $v\rho^2$ gives
    $q^{8\ell+2m}v^{2m+2}$.
    Increasing both exponents to $8\ell+2m+2$ proves
    (\ref{eq:selected_common_natural_cardinality}).
\end{proof}

\begin{theorem}[The low-frequency numerical consequence]
    \label{thm:selected_common_natural_low_frequency_cardinality}
    \lean{TNLean.MPS.AreaLaw.actual_common_physical_mpo_card_low_frequency_le}
    Let $q\ge2$, $k\ge1$, $\ell\ge1$, and let $i,j,n$ be nonnegative integers
    with $i,j\le\lfloor k/(2\ell)\rfloor$. Then
    \begin{align}
        n\le\bigl(q(k+1)\bigr)^{8\ell+2(i+j)+2}
        \quad\Longrightarrow\quad
        n\le(qk)^{20\ell+4\lfloor k/\ell\rfloor}.
        \label{eq:selected_common_natural_low_frequency_cardinality}
    \end{align}
    This is an arithmetic implication and supplies no further operator or
    bond assertion.
\end{theorem}
\begin{proof}
    The frequency bounds give $(i+j)\ell\le k$, so
    $i+j\le\lfloor k/\ell\rfloor$. Since $\ell\ge1$,
    $8\ell+2(i+j)+2\le10\ell+2\lfloor k/\ell\rfloor$.
    Moreover $q(k+1)\le(qk)^2$. Substituting this base comparison doubles the
    last exponent and gives the stated bound.
\end{proof}

\section{A common middle family after a change of interval}

For a finite set of sites $S$, write $\mathcal C(S)=\{0,\ldots,q-1\}^S$.
Restriction of configurations identifies the full chain with the tensor
product of a region and its complement. In the identities below the
complement factor is listed first. Open matrix product representations
use the paired row and column alphabet, in increasing site order.


\begin{theorem}[Natural coordinates of one transferred interval factor]
    \label{thm:interval_transfer_entry_coordinates}
    \lean{TNLean.MPS.AreaLaw.exists_intervalSlice_changingRegionalMiddleFactor}
    \uses{def:regional_middle_transferred_family,def:physical_interval_increasing_coordinates}
    Let $q\ge1$ and let $E=\lbrack c,d)$ and $J=\lbrack u,v)$ in a $T$-site chain satisfy
    $c\le d\le T$, $u\le v\le T$ and $\max(c,u)\le\min(d,v)$.
    Supply one matrix $V$ on $\mathcal C(E)$ and fixed configurations on
    $E\setminus J$ and $J\setminus E$ defining the transferred matrix $W$.
    In increasing coordinates, there are nonnegative integers $a,b,l$ with
    $a+l\le d-c$ and $b+l\le v-u$, old full configurations $f,g$, and new
    outside configurations $\eta,\theta$ such that
    \begin{align}
        W(\sigma,\tau)
         & =\mathbf1_{\sigma|_{[b,b+l)^c}=\eta}
        \mathbf1_{\tau|_{[b,b+l)^c}=\theta}
        V\bigl(\iota_a(f,\sigma|_{[b,b+l)}),
        \iota_a(g,\tau|_{[b,b+l)})\bigr).
        \label{eq:interval_transfer_entry_coordinates}
    \end{align}
    Here $V,W$ are read in their natural coordinates, and $\iota_a$ has
    the meaning in Theorem~\ref{thm:trace_interval_slice_configuration_extend}.
    Complements in this formula are in the new natural interval.
    The overlap length $l$ may be zero. No matrix product representation
    is assumed.
\end{theorem}
\begin{proof}
    Take the increasing coordinates of $E\cap J$. Extend the supplied old
    and new outside configurations to full natural configurations, using
    $q\ge1$ for unused entries. Evaluation of the prescribed regional
    matrix gives (\ref{eq:interval_transfer_entry_coordinates}).
\end{proof}

\begin{theorem}[Transfer of a common middle family between intervals]
    \label{thm:interval_mpo_common_family_transfer}
    \lean{TNLean.MPS.AreaLaw.changingRegionalMiddleFamily_mpo_and_global_decomposition}
    \leanok
    \uses{def:regional_middle_transferred_family,
        thm:regional_middle_exterior_transfer,
        def:physical_interval_increasing_coordinates,
        thm:physical_interval_subinterval_coordinates}
    Let $q$ be a positive integer, and let $c,d,u,v,T$ be nonnegative
    integers satisfying
    \begin{align}
        c\le d\le T,\qquad u\le v\le T,\qquad
        \max(c,u)\le\min(d,v).
        \label{eq:interval_mpo_transfer_endpoints}
    \end{align}
    Put $E={[c,d)}$,
    $J={[u,v)}$, $A=E\setminus J$, $B=J\setminus E$, and
    $C=E\cap J$. Let $\iota$ be a finite index set and $D:\iota\to\N$.
    Suppose that for each
    $t\in\iota$, the matrix $V_t$ on $\mathcal C(E)$ has an open matrix
    product representation with internal bond dimensions at most $D_t$.

    Define the prescribed new family, indexed by
    \begin{align}
        I=\iota\times\mathcal C(A)^2\times\mathcal C(B)^2,
        \label{eq:interval_mpo_transfer_index}
    \end{align}
    by the entries
    \begin{align}
         & W_{t,\alpha_r,\alpha_c,\eta_r,\eta_c}(\sigma,\tau)\notag   \\
         & \quad=\mathbf1_{\sigma|_B=\eta_r}\mathbf1_{\tau|_B=\eta_c}
        V_t\bigl(\alpha_r\sqcup\sigma|_C,
        \alpha_c\sqcup\tau|_C\bigr).
        \label{eq:interval_mpo_transfer_factor}
    \end{align}
    Here $\sigma,\tau\in\mathcal C(J)$, and the disjoint union combines
    configurations on $A$ and $C$ into a configuration on $E$.
    Then
    \begin{align}
        |I|=|\iota|q^{2|A|+2|B|},
        \label{eq:interval_mpo_transfer_count}
    \end{align}
    and the member indexed by $(t,\alpha_r,\alpha_c,\eta_r,\eta_c)$ has an
    open matrix product representation on $J$, in increasing site order,
    with internal bond dimensions at most $\max(1,D_t)$.

    The family $W$ is fixed before any exterior coefficient matrices are
    supplied. For every family $M_t$ on $\mathcal C(E^c)$ there are
    matrices $M'_r$ on $\mathcal C(J^c)$ such that the following two
    expressions define the same whole-chain operator:
    \begin{align}
        X & =\sum_{t\in\iota}M_t\otimes V_t
        =\sum_{r\in I}M'_r\otimes W_r.
        \label{eq:interval_mpo_transfer_decomposition}
    \end{align}
    The two expressions use their respective canonical regional
    configuration identifications. An empty index set and an overlap of
    length zero are allowed.
    This is supporting interval-transport algebra for
    \cite[Proposition~5.7]{Arad2016RigorousRG}.
\end{theorem}
\begin{proof}\leanok
    In natural coordinates, restrict each old matrix to the interval
    $E\cap J$ by fixing its row and column configurations on $E\setminus J$.
    Extend the result to $J$ by the indicated fixed row and column
    configurations on $J\setminus E$. Interval restriction does not
    increase the open bond bound, and extension gives the bound
    $\max(1,D_t)$. Positive $q$ supplies a fixed configuration on the
    common interval when identifying these operations with
    (\ref{eq:interval_mpo_transfer_factor}). The endpoint inequalities
    give precisely the increasing coordinates of the overlap, including
    an overlap of length zero.

    The configuration dimensions are $q^{|A|}$ and $q^{|B|}$, giving
    (\ref{eq:interval_mpo_transfer_count}). Finally expand the transferred
    old-only and new-only coordinates in matrix units. Theorem~\ref{thm:regional_middle_exterior_transfer} then gives
    (\ref{eq:interval_mpo_transfer_decomposition}), with the prescribed
    middle factors and the exterior coefficients chosen afterwards.
\end{proof}

\section{A common family on the prescribed middle interval}

Write $\mathcal C(S)=\{0,\ldots,q-1\}^S$ for the configurations on a
set of sites $S$. Regional tensor products list the complement first.
The following are supporting geometric and counting results for
\cite[Proposition~5.7]{Arad2016RigorousRG}.

\begin{theorem}[Bounds for the selected and prescribed intervals]
    \label{thm:common_middle_interval_bounds}
    \lean{TNLean.MPS.AreaLaw.selected_middle_fixed_interval_bounds}
    \leanok
    Let $L,a,b,\ell$ be nonnegative integers with $a+2\ell\le b$ and
    $b+2\ell\le L-1$, and let $0\le p,s<2\ell$. Put
    $c=a+p+1$, $d=b+s+1$, $u=a+\ell$, and $v=b+\ell+1$. Then
    $c\le d\le L$, $u\le v\le L$, and $\max(c,u)<\min(d,v)$.
\end{theorem}
\begin{proof}\leanok
    The bounds $p,s<2\ell$, together with the two separation inequalities,
    give each endpoint comparison.
\end{proof}

\begin{theorem}[Low-frequency cardinality on the prescribed interval]
    \label{thm:common_middle_low_frequency_cardinality}
    \lean{TNLean.MPS.AreaLaw.actual_common_fixed_middle_mpo_card_low_frequency_le}
    \leanok
    \uses{thm:selected_common_natural_cardinality,
        thm:selected_common_natural_low_frequency_cardinality}
    Retain the geometric data and positions above. Set $E=[a+p+1,b+s+1)$,
    $Z=[a+2\ell+1,b)$, and $\rho=|\mathcal C(E\setminus Z)|$.
    Let $\mathcal W_{k,i,j}$ be the two-marker index set with $q^2$ labels
    for each marker family. Suppose $q\ge2$, $k\ge1$, and
    $i,j\le\lfloor k/(2\ell)\rfloor$. For every nonnegative integer $n$,
    \begin{align}
        n\le |\mathcal W_{k,i,j}|(k+1)\rho^2q^{4\ell}
        \quad\Longrightarrow\quad
        n\le(qk)^{24\ell+4\lfloor k/\ell\rfloor}.
        \label{eq:common_middle_low_frequency_count}
    \end{align}
\end{theorem}
\begin{proof}\leanok
    The chosen position implies $\ell\ge1$. The selected-family count and its
    low-frequency consequence give
    $|\mathcal W_{k,i,j}|(k+1)\rho^2\le(qk)^{20\ell+4\lfloor k/\ell\rfloor}$.
    Multiplication by $q^{4\ell}\le(qk)^{4\ell}$ proves
    (\ref{eq:common_middle_low_frequency_count}).
\end{proof}

\begin{theorem}[A common prescribed-middle family for finite sampling nodes]
    \label{thm:common_middle_grid_mpo_family}
    \lean{TNLean.MPS.AreaLaw.exists_common_fixed_middle_grid_mpo_family}
    \leanok
    \uses{thm:common_middle_interval_bounds,thm:selected_common_natural_mpo_family,
        thm:regional_middle_from_numbered_three_regions,
        thm:interval_mpo_common_family_transfer,
        thm:interval_trace_common_family_transfer,def:regional_middle_transferred_family}
    Retain the geometric data and positions above, and suppose $q\ge1$.
    Let $B,k,i,j$ be arbitrary nonnegative integers. Supply window operators
    $W^L_r,W^R_r$, supported on $[a+r,a+r+2)$ and
    $[b+r,b+r+2)$,
    respectively. Supply site-dependent core tensors of virtual dimension $B$
    on $[a+2\ell,b+1)$ and an arbitrary closing matrix. Let $P_M$ be their
    trace contraction, canonically placed on the full chain.

    Let $\zeta$ be a finite set of sampling nodes. Use the labels
    $L_r,R_r,B_0,B_1,B_2$ and set
    $O=\{L_r,R_r,B_0,B_1,B_2\}\setminus\{L_p,R_s\}$.
    Supply complex weights $z_{\nu,t}$ for $\nu\in\zeta$ and all labels $t$.
    Put $E=[a+p+1,b+s+1)$,
    $J=[a+\ell,b+\ell+1)$,
    $Z=[a+2\ell+1,b)$, and $\rho=|\mathcal C(E\setminus Z)|$. Define
    $Q=|\mathcal W_{k,i,j}|(k+1)\rho^2$.

    There exist nonnegative integers $n_\nu$, degrees $0\le d_{\nu,r}\le k$, and matrices
    $V_{\nu,r}$ on the natural coordinates of $E$, with $0\le r<n_\nu$.
    Read these matrices on $\mathcal C(E)$ using the increasing enumeration
    of $E$. Let $A=E\setminus J$ and $B'=J\setminus E$, and define
    \begin{align}
        I_\nu                   & =\{0,\ldots,n_\nu-1\}\times\mathcal C(A)^2
        \times\mathcal C(B')^2, &
        I                       & =\coprod_{\nu\in\zeta} I_\nu.
        \label{eq:common_middle_node_indices}
    \end{align}
    For $\sigma,\tau\in\mathcal C(J)$, prescribe the middle matrices by
    \begin{align}
        W_{\nu,r,\alpha_r,\alpha_c,\eta_r,\eta_c}(\sigma,\tau)
         & =\mathbf1_{\sigma|_{B'}=\eta_r}\mathbf1_{\tau|_{B'}=\eta_c}
        V_{\nu,r}(\alpha_r\sqcup\sigma|_{E\cap J},
        \alpha_c\sqcup\tau|_{E\cap J}).
        \label{eq:common_middle_node_matrices}
    \end{align}
    The disjoint unions in (\ref{eq:common_middle_node_matrices}) combine
    configurations on $A$ and $E\cap J$ into configurations on $E$.
    The resulting family satisfies $|I|\le|\zeta|Qq^{4\ell}$.
    Each $W_\omega$ has an open matrix product representation on $J$, on the
    paired physical alphabet in increasing site order, with internal bond
    dimensions at most $(\max(1,B))^{2k}$.
    In addition, for every $\omega\in I$ there are a nonnegative integer
    $\chi_\omega\le(\max(1,B))^k$, a closing matrix
    $C_\omega\in\C^{\chi_\omega\times\chi_\omega}$, and site-dependent matrices
    $A_{\omega,v}(x,y)\in\C^{\chi_\omega\times\chi_\omega}$ such that
    \begin{align}
        W_\omega(\sigma,\tau)
         & =\operatorname{tr}\left(C_\omega
        \prod_{v=0}^{|J|-1}A_{\omega,v}(\sigma_{j_v},\tau_{j_v})\right),
        \label{eq:common_middle_trace_presentation}
    \end{align}
    where $j_0<\cdots<j_{|J|-1}$ enumerate $J$ and
    $\sigma,\tau\in\mathcal C(J)$.
    The dimension, closing matrix and tensors are chosen for the same
    $W_\omega$ before the exterior operators and scalar coefficients.

    These choices precede every supported exterior pair $P_L,P_R$ on
    $[0,a+1)$ and
    $[b+2\ell,L)$, respectively. For such a pair, set
    $F_{L_r}=W^L_r$, $F_{R_r}=W^R_r$,
    $F_{B_0}=P_L$, $F_{B_1}=P_M$, $F_{B_2}=P_R$, and
    $D_\nu=\sum_{t\in O}z_{\nu,t}F_t$.
    For every scalar array $c_{\nu,h}$, with $\nu\in\zeta$ and $0\le h\le k$,
    there are joint exterior matrices $M_\omega$ on $\mathcal C(J^c)$ such that
    \begin{align}
        \sum_{\nu\in\zeta}\sum_{h=0}^{k}
        c_{\nu,h}[x^iy^j](D_\nu+xW^L_p+yW^R_s)^h
         & =\left(\sum_{\omega\in I}M_\omega\otimes W_\omega\right)^{\mathrm{phys}}.
        \label{eq:common_middle_node_decomposition}
    \end{align}
    The superscript uses the canonical restriction of a full configuration
    to $J^c$ and $J$; the formal variables commute with matrix coefficients.
    Empty sampling sets, zero virtual dimension, and degree zero are allowed.
\end{theorem}
\begin{proof}\leanok
    \uses{thm:selected_common_natural_mpo_family,
        thm:regional_middle_from_numbered_three_regions,
        thm:interval_mpo_common_family_transfer,
        thm:interval_trace_common_family_transfer,thm:common_middle_interval_bounds,
        thm:chain_mps_of_cut_rank,thm:chain_cut_rank_of_mps}
    Choose the selected natural family separately at each supplied node.
    Pass its joint exterior decomposition to regional coordinates, then
    transfer the middle from $E$ to $J$. The geometric inequalities give the
    required overlap. Moreover $|E\setminus J|+|J\setminus E|\le2\ell$,
    so the canonical family (\ref{eq:common_middle_node_matrices}) has at most
    $Qq^{4\ell}$ members at each node. Taking the disjoint union
    (\ref{eq:common_middle_node_indices}) gives the asserted total count.

    A member of degree $t\le k$ first has open bond bound
    $\max(1,(B^t)^2)$. Since
    $\max(1,(B^t)^2)\le(\max(1,B))^{2k}$, the internal-cut-rank criterion
    gives the uniform bound. The same member is a trace contraction of
    stacked tensors of dimension $B^t$. Restricting a trace contraction to
    fixed configurations outside a subinterval replaces its closing matrix
    $C$ by $SCP$, where $P$ and $S$ are the ordered products before and after
    the subinterval. Inserting fixed row and column configurations at the
    additional sites uses the letters
    $\mathbf1_{x=\eta_v}\mathbf1_{y=\theta_v}\Id_{B^t}$.
    Thus the prescribed transfer leaves the virtual dimension $B^t$ unchanged,
    and $B^t\le(\max(1,B))^k$ gives
    (\ref{eq:common_middle_trace_presentation}).
    These operations apply to the matrices already prescribed in
    (\ref{eq:common_middle_node_matrices}).

    Finally choose the exterior matrices at each
    node for the scalar family $c_{\nu,h}$ and sum their decompositions.
    The finite sum over the disjoint union is exactly
    (\ref{eq:common_middle_node_decomposition}).
\end{proof}

\section{Common scalar recovery of two marked coefficients}

These are the interpolation and ordered-word identities used in
\cite[Lemma~4.2]{Arad2013AreaLaw} and
\cite[Proposition~5.7]{Arad2016RigorousRG}.
Only a spanning family of evaluations is required; the coefficient operators
need not be linearly independent. The distinction between spanning and a
basis in the source interpolation argument is recorded in
\cite{gap:aklv13_interpolation_and_word_count}.

\begin{definition}[Ordinary labels]
    \label{def:two_marker_ordinary_labels}
    \lean{MvPolynomial.twoMarkerOrdinaryLabels}
    For a finite alphabet $I$ and distinguished labels $u,v\in I$, let
    \begin{align}
        \mathcal O_{u,v}=I\setminus\{u,v\}.
        \label{eq:two_marker_ordinary_labels}
    \end{align}
    The definition does not require the two labels to be distinct.
\end{definition}

\begin{definition}[The full-frequency ordered-word coefficient]
    \label{def:full_frequency_word_coefficient}
    \lean{WordMultiplicity.wordCoefficient}
    Let $I$ be finite, let $R$ be a semiring, and let $H_a\in R$ for $a\in I$.
    For $r\ge0$ and a count vector $\nu:I\to\N$, put
    \begin{align}
        \mathcal W_{H,r}(\nu)=
        \sum_{\substack{w:\{0,\ldots,r-1\}\to I \\
                |w^{-1}(a)|=\nu(a)\ (a\in I)}}
        H_{w(0)}H_{w(1)}\cdots H_{w(r-1)}.
        \label{eq:full_frequency_word_coefficient}
    \end{align}
    The factors retain their displayed order. An empty product is $1$, and a
    sum with no words is zero.
\end{definition}

\begin{theorem}[One sample family for every degree and frequency slice]
    \label{thm:two_window_common_scalar_recovery}
    \lean{Polynomial.exists_twoWindow_wordCoefficient_recovery}
    \uses{def:two_marker_ordinary_labels,def:full_frequency_word_coefficient}
    Let $\ell,K$ be nonnegative integers, let
    \begin{align}
        I=(\{0,\ldots,2\ell-1\}\sqcup\{0,\ldots,2\ell-1\})
        \sqcup\{0,1,2\},
        \label{eq:two_window_scalar_alphabet}
    \end{align}
    and choose $0\le p,s<2\ell$.
    Let $u$ be the label $p$ in the first window and $v$ the label $s$ in the
    second window. Let $R$ be a ring and a complex algebra, without any
    commutativity assumption on $R$, and put $N=(K+1)^{4\ell+1}$.

    There are complex sample points $z_t:I\to\C$, $0\le t<N$, with the
    following property. For every $i,j\ge0$ and every scalar mask
    $\mu:(I\to\N)\to\C$, there are complex coefficients $\gamma_t$ such
    that, for every family $H:I\to R$ and every $0\le r\le K$,
    \begin{align}
         & \sum_{\substack{\nu:I\to\N         \\ \sum_a\nu(a)=r\\
                \nu(u)=i,\ \nu(v)=j}}
        \mu(\nu)\,\mathcal W_{H,r}(\nu)\notag \\
         & \qquad=\sum_{t=0}^{N-1}\gamma_t
        [x^iy^j]\left(D_t+yH_v+xH_u\right)^r,
        \qquad D_t=\sum_{a\in\mathcal O_{u,v}}z_t(a)H_a.
        \label{eq:two_window_common_scalar_recovery}
    \end{align}
    Here $x,y$ are central formal variables and scalar products mean the
    complex algebra action. The selected terms $H_u,H_v$ remain unweighted.
    The points precede $i,j,\mu$, while the recovery coefficients precede $H$
    and every degree $r\le K$. Points for different selected pairs need not
    coincide. The case $K=0$ is included.
\end{theorem}
\begin{proof}
    \uses{def:two_marker_ordinary_labels,def:full_frequency_word_coefficient}
    The two distinguished labels are distinct, so the ordinary label set has
    $4\ell+1$ members. Ordinary count vectors with each coordinate between
    zero and $K$ form a monomial box of exactly $(K+1)^{4\ell+1}$ elements.
    Choose the finite interpolation points for this box before the frequencies
    or mask are supplied.

    Fix $i,j,\mu$. Restore the two distinguished counts in each ordinary box
    exponent, and use $\mu$ of that restored count vector as its scalar weight.
    Uniform scalar interpolation gives coefficients $\gamma_t$ before any
    module-valued coefficient family is supplied. A frequency slice at any
    $r\le K$ embeds in the same box; set the other coefficients to zero.
    Every count coordinate in that slice is at most $K$, so restoring and
    removing the distinguished counts are inverse on the retained slice.

    Apply this recovery identity to the coefficients
    $\mathcal W_{H,r}(\nu)$. The existing counted-word identity identifies
    each evaluated ordinary monomial sum with
    $[x^iy^j](D_t+yH_v+xH_u)^r$, without changing the order of the factors
    $H_a$. Reindexing the monomial box by its exact cardinality gives
    (\ref{eq:two_window_common_scalar_recovery}).
\end{proof}

\section{A common recovery identity for two-window polynomials}

The following combines the low-frequency partition and scalar interpolation
in \cite[Lemma~4.2]{Arad2013AreaLaw} and
\cite[Proposition~5.7]{Arad2016RigorousRG}.
Interpolation uses a spanning family of evaluations, as explained in
\cite{gap:aklv13_interpolation_and_word_count}.

\begin{definition}[Crossing positions and low frequencies]
    \label{def:crossing_frequency_index}
    \lean{TNLean.MPS.AreaLaw.crossingFrequencyIndex}
    \leanok
    For $k\ge0$, set $\mathcal I_{0,k}=\varnothing$. For $\ell\ge1$, put
    $m=\lfloor k/(2\ell)\rfloor$ and
    \begin{align}
        \mathcal I_{\ell,k}
         & =\{0,\ldots,2\ell-1\}^2\times\{0,\ldots,m\}^2.
        \label{eq:crossing_frequency_index}
    \end{align}
    An index $\nu=(p,s,i,j)$ records one position in each window and their two
    frequencies.
\end{definition}

\begin{theorem}[Common recovery of every polynomial through a fixed degree]
    \label{thm:global_two_window_polynomial_recovery}
    \lean{TNLean.MPS.AreaLaw.exists_global_two_window_polynomial_recovery}
    \leanok
    \uses{def:crossing_frequency_index,thm:two_window_common_scalar_recovery,
        def:two_marker_coefficient}
    Let $\ell\ge1$ and $k\ge0$. Let $R$ be a ring and a complex algebra,
    without a commutativity assumption. Set
    $\Lambda=(\{0,\ldots,2\ell-1\}\sqcup\{0,\ldots,2\ell-1\})
        \sqcup\{0,1,2\}$ and $N=(k+1)^{4\ell+1}$.
    Write $L_p$ and $R_s$ for the two window labels.
    There are complex scalar arrays
    $z_{\nu,t}:\Lambda\to\C$ and $\gamma_{\nu,t}\in\C$, with
    $\nu=(p,s,i,j)\in\mathcal I_{\ell,k}$ and $0\le t<N$, such that the
    following identity holds for every family $H:\Lambda\to R$ and all
    complex coefficients $c_0,\ldots,c_k$. Define
    \begin{align}
        D_{\nu,t} & =\sum_{a\in\Lambda\setminus\{L_p,R_s\}}
        z_{\nu,t}(a)H_a.
        \label{eq:global_two_window_ordinary_sum}
    \end{align}
    Then
    \begin{align}
        \sum_{r=0}^{k} c_r\left(\sum_{a\in\Lambda}H_a\right)^r
         & =\sum_{\nu=(p,s,i,j)\in\mathcal I_{\ell,k}}
        \sum_{t=0}^{N-1}\sum_{r=0}^{k}
        c_r\gamma_{\nu,t}\,[x^iy^j]
        (D_{\nu,t}+xH_{L_p}+yH_{R_s})^r.
        \label{eq:global_two_window_recovery}
    \end{align}
    The variables $x,y$ are central, and scalar multiplication is the complex
    algebra action. The selected terms are unweighted. Both arrays are fixed
    before $H$ and the entire polynomial coefficient family are supplied.
    The case $k=0$ is included. This is a scalar recovery identity; it imposes
    no spatial support, spectral, tensor rank, or matrix product hypothesis.
\end{theorem}
\begin{proof}\leanok
    \uses{def:crossing_frequency_index,thm:two_window_common_scalar_recovery,
        def:full_frequency_word_coefficient,def:two_marker_coefficient}
    The common-cutoff word partition chooses one selector on the frequency
    vectors of each window. If a full frequency vector has total at most $k$,
    the selected count in each window is at most
    $\lfloor k/(2\ell)\rfloor$. These same selectors work before $H$ and at
    every degree $r\le k$.

    For $\nu=(p,s,i,j)$, use the scalar mask equal to one when the two selectors
    choose $L_p,R_s$, and zero otherwise. The words assigned to this bin are
    exactly the full-frequency slice with counts $i,j$, weighted by this mask.
    The two-window scalar recovery theorem supplies $z_{\nu,t}$ and
    $\gamma_{\nu,t}$ before $H$ and all degrees. These masks depend only on
    the selectors and the index $\nu$.

    Partition each power into these bins and replace each masked slice by its
    recovery identity. Reindex the window labels and their bounded counts by
    (\ref{eq:crossing_frequency_index}). Multiplying by $c_r$ and interchanging
    the finite degree and sample sums gives
    (\ref{eq:global_two_window_recovery}), with ordinary sums
    (\ref{eq:global_two_window_ordinary_sum}). All ordered products are preserved.
\end{proof}

\section{A common physical middle family for every polynomial}

Let the chain have $L$ sites and local space $\C^q$.
For a region $S$, write $\mathcal C(S)=\{0,\ldots,q-1\}^S$.
Every regional tensor product below lists the complement before the middle.
All intervals are intersected with the chain.

Fix nonnegative integers $a,b,\ell,B,k$ satisfying
\begin{align}
    a+2\ell\le b,\qquad b+2\ell\le L-1,\qquad
    \ell\ge1,\quad q\ge2,\quad k\ge1.
    \label{eq:physical_polynomial_family_hypotheses}
\end{align}
Supply full-chain window operators $W^L_r,W^R_r$, $0\le r<2\ell$,
supported on $[a+r,a+r+2)$ and
$[b+r,b+r+2)$, respectively.
Set $K=[a+2\ell,b+1)$ and supply site-dependent tensors
$A_t(i,j)\in\C^{B\times B}$ for $0\le t<|K|$, together with an arbitrary
closing matrix $C_0\in\C^{B\times B}$. Let $P_M$ be the canonical
full-chain extension of their trace contraction
\begin{align}
    (\sigma,\tau)\longmapsto
    \operatorname{tr}\left(C_0\prod_{t=0}^{|K|-1}A_t(\sigma_t,\tau_t)\right),
    \label{eq:physical_polynomial_supplied_core}
\end{align}
using the increasing enumeration of $K$. Put
\begin{align}
    J          & =[a+\ell,b+\ell+1),                    &
    D          & =(qk)^{18\ell+2\lfloor k/\ell\rfloor}, &
    \mathcal B & =(\max(1,B))^{2k}.
    \label{eq:physical_polynomial_two_bounds}
\end{align}
The count parameter $D$ and the open-boundary bond bound $\mathcal B$
have distinct roles.

These are supporting polynomial-family results for
\cite[Proposition~5.7]{Arad2016RigorousRG}.
The interpolation argument requires only spanning, without independence of
the coefficient operators; see
\cite{gap:aklv13_interpolation_and_word_count}.

\begin{theorem}[The canonical common middle family]
    \label{thm:physical_polynomial_canonical_middle_family}
    \lean{TNLean.MPS.AreaLaw.exists_common_physical_polynomial_mpo_family}
    \leanok
    \uses{def:crossing_frequency_index,thm:global_two_window_polynomial_recovery,
        thm:common_middle_grid_mpo_family,
        thm:common_middle_low_frequency_cardinality,
        def:regional_middle_transferred_family,thm:chain_two_window_family_cardinality}
    Let
    \begin{align}
        \kappa=\{0,\ldots,2\ell-1\}^2\times
        \{0,\ldots,\lfloor k/(2\ell)\rfloor\}^2,
        \qquad m=(k+1)^{4\ell+1}.
        \label{eq:physical_polynomial_crossing_index}
    \end{align}
    For $v=(p,s,i,j)\in\kappa$, put
    $E_v=[a+p+1,b+s+1)$.
    There are nonnegative integers $n_{v,t}$ and matrices $V_{v,t,r}$ in the
    increasing natural coordinates of $E_v$, for $0\le t<m$ and
    $0\le r<n_{v,t}$, which determine the following prescribed family.
    Write $V^{\mathrm{reg}}_{v,t,r}$ for these matrices read on
    $\mathcal C(E_v)$, and set
    \begin{align}
        I_{v,t} & =\{0,\ldots,n_{v,t}-1\}\times
        \mathcal C(E_v\setminus J)^2\times\mathcal C(J\setminus E_v)^2,
                & I                             & =\coprod_{v\in\kappa}\coprod_{t=0}^{m-1} I_{v,t}.
        \label{eq:physical_polynomial_canonical_index}
    \end{align}
    Apply the prescribed regional middle-family transfer to each family
    $V^{\mathrm{reg}}_{v,t}$ from $E_v$ to $J$, and denote the resulting
    matrices by $W_\omega$, $\omega\in I$. Then $|I|\le D^2$, and each
    $W_\omega$, on the paired physical alphabet in increasing site order on
    $J$, has an open-boundary matrix product representation with internal
    bond dimensions at most $\mathcal B$.

    All these choices precede the following data. Supply arbitrary full-chain
    operators $P_L,P_R$ supported on $[0,a+1)$ and
    $[b+2\ell,L)$, and put
    \begin{align}
        H=P_L+P_M+P_R+\sum_{r=0}^{2\ell-1}(W^L_r+W^R_r).
        \label{eq:physical_polynomial_hamiltonian}
    \end{align}
    For every scalar family $c_0,\ldots,c_k\in\C$, there are matrices
    $M_\omega$ on $\mathcal C(J^c)$ such that
    \begin{align}
        \sum_{r=0}^{k}c_rH^r
        =\left(\sum_{\omega\in I}M_\omega\otimes W_\omega\right)^{\mathrm{phys}}.
        \label{eq:physical_polynomial_canonical_decomposition}
    \end{align}
    The superscript denotes the canonical restriction of a full configuration
    to $J^c$ and $J$. No self-adjointness, positivity or gap assumption is made.
    Zero virtual dimension is allowed. This is the general square-derived
    bond bound, with no sharper source bond or running-time assertion.
\end{theorem}
\begin{proof}\leanok
    \uses{thm:global_two_window_polynomial_recovery,thm:common_middle_grid_mpo_family,
        thm:common_middle_low_frequency_cardinality,
        thm:chain_two_window_family_cardinality}
    Choose the common scalar recovery points and weights before the operator
    family and the polynomial coefficients. For every crossing position and
    low-frequency pair, apply the supplied-node middle-family theorem at those
    points. Take the dependent union of its canonical carriers, retaining each
    prescribed regional middle operator. The low-frequency estimates bound the
    number of terms before the final numerical absorption. There are
    $(k+1)^{4\ell+1}\le(k+1)^{4\ell+3}$ scalar nodes, so the square
    cardinality estimate gives $|I|\le D^2$.

    After $P_L,P_R$ and the polynomial coefficients are supplied, scalar
    recovery expresses the polynomial as a sum of the corresponding
    two-frequency coefficients. Choose their exterior matrices in the common
    family and sum them, giving
    (\ref{eq:physical_polynomial_canonical_decomposition}).
\end{proof}

\begin{theorem}[A finite enumeration of the common middle family]
    \label{thm:physical_polynomial_fin_middle_family}
    \lean{TNLean.MPS.AreaLaw.exists_common_physical_polynomial_mpo_family_fin}
    \leanok
    \uses{thm:physical_polynomial_canonical_middle_family}
    Under the same hypotheses and with the same supplied windows, core tensors
    and cap, there are a nonnegative integer $N\le D^2$ and matrices
    $W_0,\ldots,W_{N-1}$ on $\mathcal C(J)$, each with natural open-boundary
    bond bound $\mathcal B$, such that for every supported $P_L,P_R$ and every
    complex family $c_0,\ldots,c_k$ there are exterior matrices $M_r$ with
    \begin{align}
        \sum_{h=0}^{k}c_hH^h
        =\left(\sum_{r=0}^{N-1}M_r\otimes W_r\right)^{\mathrm{phys}}.
        \label{eq:physical_polynomial_fin_decomposition}
    \end{align}
    The number and middle matrices are chosen before both the exterior pair and
    all polynomial coefficients.
\end{theorem}
\begin{proof}\leanok
    \uses{thm:physical_polynomial_canonical_middle_family}
    Enumerate the finite canonical carrier $I$ by
    $\{0,\ldots,|I|-1\}$. Pull both the middle matrices and the later exterior
    matrices back along this bijection. The cardinality, bond bounds and finite
    sum identity are unchanged.
\end{proof}

\begin{theorem}[A common middle family for polynomial evaluation]
    \label{thm:physical_aeval_fin_middle_family}
    \lean{TNLean.MPS.AreaLaw.exists_common_physical_aeval_mpo_family_fin}
    \leanok
    \uses{thm:physical_polynomial_fin_middle_family}
    Under the same hypotheses and fixed data, one may choose $N\le D^2$ and
    middle matrices $W_0,\ldots,W_{N-1}$ with the same natural bond bound
    $\mathcal B$ before every supported exterior pair and every polynomial
    $P\in\C[X]$ of natural degree at most $k$, so that suitable exterior
    matrices satisfy
    \begin{align}
        P(H)=\left(\sum_{r=0}^{N-1}M_r\otimes W_r\right)^{\mathrm{phys}}.
        \label{eq:physical_aeval_fin_decomposition}
    \end{align}
    The zero polynomial is included, with no normalization condition on $P$.
\end{theorem}
\begin{proof}\leanok
    \uses{thm:physical_polynomial_fin_middle_family}
    Apply the preceding theorem with $c_h$ equal to the coefficient of $X^h$
    in $P$. The finite polynomial-evaluation identity gives
    $P(H)=\sum_{h=0}^{k}c_hH^h$. It remains valid in the noncommutative matrix
    algebra because the complex scalar coefficients act centrally.
\end{proof}

\subsection{Changing the middle tensor factor}

The change of cuts in~\cite[Proposition~5.7]{Arad2016RigorousRG}
uses only matrix units to transfer the sites removed from, or added to, the
middle interval. The following identities isolate this finite-dimensional
algebra; they impose no ordering or bond bound on the matrices.

\begin{definition}[Exchange of the middle coordinates]
    \label{def:chain_middle_factor_exchange}
    \lean{Matrix.middleFactorExchangeEquiv,Matrix.exchangedMiddleFactor,Matrix.exchangedMiddleFamily}
    \leanok
    Let $A,B,C,O$ be arbitrary sets. Define the bijection
    \begin{align}
        \Phi\colon (B\times O)\times(A\times C)
                      & \longrightarrow(A\times O)\times(B\times C),\notag \\
        ((b,o),(a,c)) & \longmapsto((a,o),(b,c)).\notag
    \end{align}
    Let $R$ be a set with distinguished elements $0,1$ and a binary
    multiplication. For arbitrary sets $Y,Z$ and a matrix $X$ on $Y\times Z$
    over $R$, write $X^{\mathrm R}_{yy'}$ for the matrix on $Z$ with entries
    $(X^{\mathrm R}_{yy'})_{z,z'}=X_{(y,z),(y',z')}$.
    Suppose equality on $B$ is decidable, and write $E_{bb'}$ for the
    standard matrix unit on $B$. For $V$ on $A\times C$, set
    \begin{align}
        V_{a,a',b,b'}=E_{bb'}\otimes V^{\mathrm R}_{aa'}.\notag
    \end{align}
    For an arbitrary index set $I$ and a family $(V_t)_{t\in I}$, apply
    this definition to each $V_t$.
    The resulting family depends only on the original middle matrices.
\end{definition}

\begin{lemma}[Matrix-unit expansion after exchanging the middle]
    \label{thm:chain_middle_factor_exchange}
    \lean{Matrix.reindex_kronecker_middleFactorExchange,Matrix.reindex_sum_kronecker_middleFactorExchange}
    \leanok
    \uses{def:chain_middle_factor_exchange}
    Let $A,B$ be finite sets with decidable equality, let $C,O$ be arbitrary
    sets, and let $R$ be a semiring. Let $M$ be a matrix on $B\times O$ and $V$
    a matrix on $A\times C$, both over $R$. Transporting both indices by $\Phi$ gives
    \begin{align}
        (M\otimes V)^\Phi
         & =\sum_{a,a'\in A}\sum_{b,b'\in B}
        \bigl(E_{aa'}\otimes M^{\mathrm R}_{bb'}\bigr)\otimes V_{a,a',b,b'}.\notag
    \end{align}
    For a finite index set $I$ and prescribed matrices $V_t$, the same
    family $(V_{t,a,a',b,b'})$ therefore gives such a decomposition of
    $(\sum_{t\in I}M_t\otimes V_t)^\Phi$ for every family $(M_t)$.
    No commutativity of $R$ or nonemptiness of the index sets is required.
\end{lemma}
\begin{proof}\leanok
    At row $((a,o),(b,c))$ and column $((a',o'),(b',c'))$, the matrix
    units leave exactly the summand with these four transferred labels.
    Its entry is
    \begin{align}
        M_{(b,o),(b',o')}\,V_{(a,c),(a',c')},\notag
    \end{align}
    which is the corresponding entry on the left, in the same
    multiplication order. This proves the first identity. Summing it over
    $t$ proves the finite-family assertion.
\end{proof}

\begin{lemma}[Number of exchanged middle factors]
    \label{thm:chain_middle_factor_exchange_card}
    \lean{Matrix.card_exchangedMiddleFamily_index}
    \leanok
    \uses{def:chain_middle_factor_exchange}
    If $I,A,B$ are finite, the prescribed family has $|I|\,|A|^2|B|^2$ indexed
    terms, including repetitions and zero matrices.
\end{lemma}
\begin{proof}\leanok
    Its index set is $I\times(A\times A)\times(B\times B)$; take cardinalities.
\end{proof}

